% =====================================================================
%  DOCUMENTACIÓN DE INVESTIGACIÓN Y PLANIFICACIÓN
%  Departamento de Informática y Computación
%  Compilar con: pdflatex -> pdflatex -> pdflatex
%  Probado en TeXStudio con MiKTeX (pdfLaTeX)
% =====================================================================
\documentclass[12pt,letterpaper,oneside]{report}

% ------------------------- Codificación e idioma --------------------
\usepackage[utf8]{inputenc}
\usepackage[T1]{fontenc}
\usepackage[spanish,es-tabla,es-nodecimaldot,es-noshorthands]{babel}
\usepackage{lmodern}
\usepackage{microtype}

% ------------------------- Geometría y diseño -----------------------
\usepackage[letterpaper,top=3cm,bottom=3cm,left=3cm,right=2.5cm,headheight=15pt]{geometry}
\usepackage{fancyhdr}
\usepackage{titlesec}
\usepackage{setspace}
\usepackage{parskip}
\onehalfspacing

% ------------------------- Gráficos y diagramas ---------------------
\usepackage{tikz}
\usepackage{pgf-umlcd}
\usepackage{pgfgantt}
\usepackage{float}
\usepackage{graphicx}
\usetikzlibrary{arrows.meta,positioning,calc,fit,backgrounds,
                decorations.pathreplacing,chains,
                shapes.geometric,shapes.misc,matrix}

% ------------------------- Tablas -----------------------------------
\usepackage{booktabs}
\usepackage{longtable}
\usepackage{array}
\usepackage{tabularx}
\usepackage{multirow}
\usepackage{colortbl}

% ------------------------- Cajas de color ---------------------------
\usepackage[most]{tcolorbox}

% ------------------------- Listas y código --------------------------
\usepackage{enumitem}
\usepackage{listings}
\usepackage{xcolor}

% ------------------------- Utilidades -------------------------------
\usepackage{hyperref}
\usepackage{url}
\usepackage{amsmath,amssymb}
\usepackage{caption}
\usepackage{subcaption}
\usepackage[nohyperlinks]{acronym}

% ------------------------- Colores institucionales ------------------
\definecolor{azulUC}{RGB}{21,101,192}
\definecolor{azulOscuro}{RGB}{13,71,161}
\definecolor{azulClaro}{RGB}{227,242,253}
\definecolor{verdeUC}{RGB}{56,142,60}
\definecolor{verdeClaro}{RGB}{232,245,233}
\definecolor{naranjaUC}{RGB}{239,108,0}
\definecolor{naranjaClaro}{RGB}{255,243,224}
\definecolor{moradoUC}{RGB}{173,20,87}
\definecolor{moradoClaro}{RGB}{243,229,245}
\definecolor{grisUC}{RGB}{96,96,96}
\definecolor{grisClaro}{RGB}{245,245,245}

% ------------------------- Hipervínculos ----------------------------
\hypersetup{
    colorlinks=true,
    linkcolor=azulOscuro,
    urlcolor=azulUC,
    citecolor=verdeUC,
    pdftitle={Documentacion de Investigacion - Departamento de Informatica y Computacion},
    pdfauthor={Departamento de Informatica y Computacion},
    pdfsubject={Metodologia de Investigacion, UML y Planificacion},
    bookmarksnumbered=true,
    bookmarksopen=true
}

% ------------------------- Encabezados y pies -----------------------
\pagestyle{fancy}
\fancyhf{}
\fancyhead[L]{\small\textcolor{azulOscuro}{\textit{Departamento de Informática y Computación}}}
\fancyhead[R]{\small\textcolor{azulOscuro}{\thepage}}
\fancyfoot[C]{\small\textcolor{grisUC}{Documentación de Investigación y Planificación}}
\renewcommand{\headrulewidth}{0.6pt}
\renewcommand{\footrulewidth}{0.4pt}

% ------------------------- Estilo de capítulos ----------------------
\titleformat{\chapter}[display]
  {\normalfont\huge\bfseries\color{azulOscuro}}
  {\filright\Large\textcolor{azulUC}{\chaptertitlename\ \thechapter}}
  {12pt}
  {\Huge}
  [\vspace{4pt}{\color{azulUC}\titlerule[1.5pt]}]

\titleformat{\section}
  {\normalfont\Large\bfseries\color{azulOscuro}}{\thesection}{1em}{}

\titleformat{\subsection}
  {\normalfont\large\bfseries\color{azulUC}}{\thesubsection}{1em}{}

\titlespacing*{\chapter}{0pt}{-20pt}{30pt}

% ------------------------- Captions --------------------------------
\captionsetup{
  font=small,
  labelfont={bf,color=azulOscuro},
  labelsep=colon,
  justification=centering
}

% ------------------------- Estilo de código -------------------------
\lstdefinestyle{codigo}{
  backgroundcolor=\color{grisClaro},
  basicstyle=\ttfamily\footnotesize,
  keywordstyle=\color{azulOscuro}\bfseries,
  commentstyle=\color{verdeUC}\itshape,
  stringstyle=\color{moradoUC},
  numberstyle=\tiny\color{grisUC},
  numbers=left,
  numbersep=8pt,
  frame=single,
  rulecolor=\color{azulUC},
  breaklines=true,
  showstringspaces=false,
  tabsize=2,
  captionpos=b
}
\lstset{style=codigo}

% ------------------------- Cajas personalizadas ---------------------
\newtcolorbox{cajaObjetivo}[1][]{
  colback=verdeClaro, colframe=verdeUC, boxrule=0.8pt, arc=3pt,
  left=6pt, right=6pt, top=6pt, bottom=6pt,
  fonttitle=\bfseries\color{white}, coltitle=white,
  colbacktitle=verdeUC, title={#1}
}

\newtcolorbox{cajaProblema}[1][]{
  colback=naranjaClaro, colframe=naranjaUC, boxrule=0.8pt, arc=3pt,
  left=6pt, right=6pt, top=6pt, bottom=6pt,
  fonttitle=\bfseries\color{white}, coltitle=white,
  colbacktitle=naranjaUC, title={#1}
}

\newtcolorbox{cajaNota}[1][]{
  colback=azulClaro, colframe=azulUC, boxrule=0.8pt, arc=3pt,
  left=6pt, right=6pt, top=6pt, bottom=6pt,
  fonttitle=\bfseries\color{white}, coltitle=white,
  colbacktitle=azulUC, title={#1}
}

\newtcolorbox{cajaAlgoritmo}[1][]{
  colback=moradoClaro, colframe=moradoUC, boxrule=0.8pt, arc=3pt,
  left=6pt, right=6pt, top=6pt, bottom=6pt,
  fonttitle=\bfseries\color{white}, coltitle=white,
  colbacktitle=moradoUC, title={#1}
}

% ------------------------- Estilos TikZ UML ------------------------
\tikzset{
  actor/.style={
    draw=azulOscuro, fill=naranjaClaro, thick, rounded corners=2pt,
    text=azulOscuro, align=center, font=\small\bfseries,
    minimum width=2.8cm, minimum height=1cm, inner sep=4pt
  },
  caso/.style={
    draw=moradoUC, fill=moradoClaro, thick, ellipse, aspect=2.6,
    text=moradoUC!80!black, align=center, font=\scriptsize,
    minimum width=3.6cm, minimum height=0.85cm, inner sep=2pt
  },
  paquete/.style={
    draw=verdeUC, fill=verdeClaro, thick, rounded corners=4pt,
    inner sep=10pt
  },
  sistema/.style={
    draw=azulUC, fill=azulClaro!40, thick, rounded corners=6pt,
    inner sep=14pt
  },
  clase/.style={
    draw=azulOscuro, fill=azulClaro, thick, rounded corners=3pt,
    text=azulOscuro, align=center, font=\small,
    minimum width=3.2cm, minimum height=1cm, inner sep=4pt
  },
  claseabs/.style={
    draw=verdeUC, fill=verdeClaro, thick, dashed, rounded corners=3pt,
    text=verdeUC!60!black, align=center, font=\small\itshape,
    minimum width=3.2cm, minimum height=1cm, inner sep=4pt
  },
  tabla/.style={
    draw=grisUC, fill=white, thick, rounded corners=2pt,
    text=black, align=left, font=\scriptsize, inner sep=4pt
  },
  flujo/.style={
    draw=azulOscuro, fill=azulClaro, thick, rounded corners=12pt,
    text=azulOscuro, align=center, font=\scriptsize\bfseries,
    minimum width=3.4cm, minimum height=0.9cm, inner sep=4pt
  },
  decision/.style={
    draw=naranjaUC, fill=naranjaClaro, thick, diamond, aspect=2.4,
    text=naranjaUC!70!black, align=center, font=\scriptsize\bfseries,
    inner sep=1pt
  },
  terminal/.style={
    draw=verdeUC, fill=verdeClaro, thick, rounded corners=14pt,
    text=verdeUC!60!black, align=center, font=\scriptsize\bfseries,
    minimum width=3cm, minimum height=0.8cm, inner sep=4pt
  },
  flecha/.style={-{Stealth[length=2mm,width=1.4mm]}, thick, color=azulOscuro},
  generalizacion/.style={-{Stealth[length=3mm,width=3mm]}, thick, color=verdeUC, dashed}
}

% Nota: pgf-umlcd reconfigura el espacio de claves /tikz, por lo que las
% relaciones punteadas de los diagramas se dibujan con la sintaxis explicita
% -{Stealth[...]} directamente en cada \draw.

% ------------------------- Comandos útiles --------------------------
\newcommand{\sistema}{\textbf{SGDIC}}
\newcommand{\sgdicshort}{SGDIC}

% =====================================================================
\begin{document}

% ============================ PORTADA ================================
\begin{titlepage}
\centering
\vspace*{0.5cm}

\begin{tikzpicture}
\node[draw=azulUC, fill=azulClaro, rounded corners=8pt, thick,
      inner sep=18pt, text width=14cm, align=center] {
  {\small\textcolor{grisUC}{UNIVERSIDAD NACIONAL DE COLOMBIA}}\\[3pt]
  {\small\textcolor{grisUC}{SEDE MANIZALES}}\\[4pt]
  {\Large\bfseries\textcolor{azulOscuro}{FACULTAD DE ADMINISTRACIÓN}}\\[4pt]
  {\large\textcolor{azulUC}{DEPARTAMENTO DE INFORMÁTICA Y COMPUTACIÓN}}
};
\end{tikzpicture}

\vspace{1.2cm}

{\Huge\bfseries\textcolor{azulOscuro}{Documentación de Investigación\\[6pt] y Planificación del Sistema}}\\[10pt]

{\large\textcolor{grisUC}{Sistema de Gestión Departamental\\ de Informática y Computación (\sistema)}}

\vspace{1cm}

\begin{tikzpicture}
\draw[azulUC, line width=2pt] (0,0) -- (13,0);
\draw[moradoUC, line width=1pt] (0,-0.12) -- (13,-0.12);
\end{tikzpicture}

\vspace{1cm}

\begin{tabular}{rl}
\textbf{\textcolor{azulOscuro}{Proyecto:}} & Sistema de Gestión Departamental \\[4pt]
\textbf{\textcolor{azulOscuro}{Área:}} & Informática y Computación \\[4pt]
\textbf{\textcolor{azulOscuro}{Versión:}} & 1.0 \\[4pt]
\textbf{\textcolor{azulOscuro}{Fecha:}} & \today \\
\end{tabular}

\vfill

{\small\textcolor{grisUC}{Documento elaborado como base metodológica para la planificación\\ y el modelado del sistema propuesto.}}
\end{titlepage}

% ======================= PÁGINAS PRELIMINARES ========================
\pagenumbering{roman}
\setcounter{page}{1}

\tableofcontents
\clearpage
\listoftables
\clearpage
\listoffigures
\clearpage

\chapter*{Resumen Ejecutivo}
\addcontentsline{toc}{chapter}{Resumen Ejecutivo}
\markboth{Resumen Ejecutivo}{Resumen Ejecutivo}

El presente documento establece la base metodológica, conceptual y de modelado
para la construcción del \textbf{Sistema de Gestión Departamental de Informática y
Computación} (\sistema). El problema identificado en el \ac{DIC} se caracteriza por
una gestión fragmentada de sus procesos académicos y administrativos: las
solicitudes de cupos, la creación y edición de materias, la comunicación entre
estudiantes y docentes, el registro de faltas, la calificación de docentes, la
gestión de quejas y sugerencias, y la verificación de resultados por parte de la
secretaría se realizan de manera aislada, manual o tardía.

La consecuencia directa es la \textbf{pérdida de oportunidad} en la toma de
decisiones: las evaluaciones docentes no se completan a tiempo, las quejas no
llegan al docente responsable en el momento adecuado, y no existe información
analítica que permita anticipar problemas de disponibilidad de horarios o
saturación de cupos.

El documento se organiza en siete capítulos: (1) introducción y planteamiento del
problema, (2) metodología de investigación, (3) misión, visión, objetivos y
valores, (4) requerimientos funcionales y no funcionales, (5) modelado \ac{UML}
mediante diagramas de casos de uso, clases, secuencia, actividades y estados,
(6) propuesta del algoritmo de soluciones rápidas con analítica de datos, y
(7) cronograma, riesgos y conclusiones.

\begin{cajaNota}[Resultados esperados del proyecto]
\begin{itemize}[leftmargin=1.2cm, itemsep=2pt]
  \item Reducción $\geq 60\%$ del tiempo de respuesta a solicitudes de cupos.
  \item \textbf{100\%} de las evaluaciones docentes completadas dentro del periodo
        académico establecido.
  \item Canal unificado de comunicación estudiante--docente--secretaría.
  \item Trazabilidad completa de quejas y sugerencias con estado y responsable.
  \item Modelo analítico que recomiende horarios óptimos según preferencias
        históricas de los estudiantes.
  \item Tablero (\textit{dashboard}) de verificación de resultados para la
        secretaría y el departamento.
\end{itemize}
\end{cajaNota}

\clearpage

\chapter*{Acrónimos y Abreviaturas}
\addcontentsline{toc}{chapter}{Acrónimos y Abreviaturas}
\markboth{Acrónimos y Abreviaturas}{Acrónimos y Abreviaturas}

\begin{acronym}[CRISP-DM]
  \acro{DIC}{Departamento de Informática y Computación}
  \acro{SGDIC}{Sistema de Gestión Departamental de Informática y Computación}
  \acro{UML}{Unified Modeling Language (Lenguaje Unificado de Modelado)}
  \acro{CRISP-DM}{Cross Industry Standard Process for Data Mining}
  \acro{RF}{Requerimiento Funcional}
  \acro{RNF}{Requerimiento No Funcional}
  \acro{KPI}{Key Performance Indicator (Indicador Clave de Desempeño)}
  \acro{API}{Application Programming Interface}
  \acro{RBAC}{Role-Based Access Control (Control de Acceso Basado en Roles)}
  \acro{ETL}{Extract, Transform, Load}
  \acro{REST}{Representational State Transfer}
  \acro{CSV}{Comma-Separated Values}
  \acro{KNN}{K-Nearest Neighbors}
  \acro{PII}{Personally Identifiable Information}
\end{acronym}

\clearpage

% ==================== CONTENIDO PRINCIPAL ===========================
\pagenumbering{arabic}
\setcounter{page}{1}

% =====================================================================
% CAPÍTULO 1: INTRODUCCIÓN Y PLANTEAMIENTO DEL PROBLEMA
% =====================================================================
\chapter{Introducción y Planteamiento del Problema}
\label{ch:introduccion}

\section{Contexto Institucional}

El \ac{DIC} es la unidad académica responsable de la formación profesional en las
áreas de informática, computación, desarrollo de software y tecnologías de la
información. Su operación diaria involucra cuatro actores fundamentales:

\begin{itemize}[leftmargin=1.4cm, itemsep=4pt]
  \item \textbf{Estudiantes:} demandan cupos, consultan horarios, califican
        docentes, reportan faltas y presentan quejas o sugerencias.
  \item \textbf{Docentes:} administran materias, registran asistencia y faltas,
        se comunican con estudiantes y responden evaluaciones.
  \item \textbf{Secretaría:} recibe, canaliza y da seguimiento a solicitudes,
        quejas y sugerencias; coordina la comunicación entre las partes.
  \item \textbf{Departamento (Dirección/Administración):} verifica resultados,
        analiza información y toma decisiones académicas.
\end{itemize}

A pesar de la naturaleza tecnológica del departamento, sus procesos
administrativos y académicos no cuentan con un sistema integrado que articule
estos actores, lo cual genera reprocesos, demoras y pérdida de información.

\section{Descripción del Problema}

\begin{cajaProblema}[Problema central]
El \ac{DIC} carece de un mecanismo unificado, trazable y analítico para gestionar
la comunicación, las solicitudes y las evaluaciones entre estudiantes, docentes,
secretaría y dirección, lo que provoca demoras en la respuesta, evaluaciones
incompletas, quejas sin seguimiento y decisiones tomadas sin evidencia.
\end{cajaProblema}

A partir del levantamiento inicial de necesidades se identificaron los siguientes
problemas específicos:

\begin{enumerate}[leftmargin=1.4cm, itemsep=5pt]
  \item \textbf{Solicitud de cupos.} El proceso se realiza de forma manual o por
        canales informales (correo, mensajería, presencial), sin control de
        disponibilidad en tiempo real ni criterios de priorización transparentes.

  \item \textbf{Creación y edición de materias.} Las materias nuevas o sus
        modificaciones (contenido, prerrequisitos, créditos, horarios) no siguen
        un flujo formal de solicitud--aprobación--publicación.

  \item \textbf{Comunicación estudiante $\rightarrow$ docente.} No existe un canal
        institucional que garantice que el mensaje llegue, sea leído y tenga
        respuesta dentro de un plazo razonable.

  \item \textbf{Comunicación docente $\rightarrow$ estudiante.} Las notificaciones
        sobre faltas, cambios de horario o novedades académicas dependen de
        medios externos y no quedan registradas.

  \item \textbf{Registro y notificación de faltas.} Las faltas se registran de
        forma tardía y el estudiante se entera cuando el acumulado ya afecta su
        situación académica.

  \item \textbf{Evaluación de docentes por estudiantes.} La evaluación no se
        aplica a tiempo, la participación es baja y los resultados se consolidan
        tarde, cuando ya no permiten retroalimentación oportuna.

  \item \textbf{Quejas y sugerencias.} Se reciben por canales dispersos, sin
        categorización, sin responsable asignado y sin trazabilidad de estado.

  \item \textbf{Efectividad de la secretaría.} La secretaría actúa como cuello de
        botella: recibe las quejas y solicitudes, pero no cuenta con herramientas
        para priorizar, delegar y comunicarse formalmente con el docente.

  \item \textbf{Verificación de resultados.} La secretaría y el departamento no
        disponen de tableros que les permitan verificar resultados, quejas y
        evaluaciones de forma consolidada.

  \item \textbf{Asignación de horarios.} Los horarios se construyen sin conocer
        las preferencias reales de los estudiantes, generando baja asistencia y
        conflictos de disponibilidad.

  \item \textbf{Ausencia de analítica.} No se explotan los datos históricos para
        anticipar saturación de cupos, docentes con bajo desempeño o patrones de
        quejas recurrentes.
\end{enumerate}

\section{Formulación del Problema}

\begin{cajaNota}[Pregunta de investigación principal]
¿Cómo puede un sistema integrado de gestión, comunicación y analítica de datos
mejorar la oportunidad, la trazabilidad y la calidad de las decisiones en el
\ac{DIC}?
\end{cajaNota}

De la pregunta principal se derivan las siguientes preguntas específicas:

\begin{itemize}[leftmargin=1.4cm, itemsep=3pt]
  \item ¿Cuáles son los procesos críticos del \ac{DIC} que presentan mayor
        demora y menor trazabilidad?
  \item ¿Qué arquitectura de información permite unificar la comunicación entre
        estudiantes, docentes y secretaría sin duplicar canales?
  \item ¿Cómo diseñar un mecanismo de evaluación docente que garantice el
        \textbf{100\%} de aplicación dentro del periodo?
  \item ¿Qué variables predicen la preferencia de horarios de los estudiantes y
        cómo pueden usarse para optimizar la programación académica?
  \item ¿Qué modelo analítico permite detectar anomalías (quejas recurrentes,
        cupos saturados, docentes con baja evaluación) y generar soluciones
        rápidas accionables?
  \item ¿Cómo verificar la efectividad de las acciones tomadas por la secretaría
        y el departamento?
\end{itemize}

\section{Justificación}

\subsection{Justificación académica}
Un departamento de informática debe ser el primer ejemplo de digitalización
institucional. Diseñar y documentar este sistema constituye un caso de estudio
aplicado sobre arquitectura de software, modelado \ac{UML} e ingeniería de datos.

\subsection{Justificación operativa}
La automatización de solicitudes de cupos, la edición de materias, el registro
de faltas y la gestión de quejas reduce el trabajo manual de la secretaría,
elimina el retrabajo y libera tiempo para actividades de mayor valor.

\subsection{Justificación estratégica}
La incorporación de analítica de datos y de un algoritmo de soluciones rápidas
convierte la operación cotidiana en información accionable: permite anticipar
problemas, priorizar intervenciones y medir el impacto de las decisiones.

\subsection{Justificación social}
Un canal formal, trazable y confidencial para quejas y sugerencias mejora la
convivencia académica, protege al estudiante y ofrece al docente
retroalimentación oportuna y justa.

\section{Objetivos}

\subsection{Objetivo general}
\begin{cajaObjetivo}[Objetivo general]
Diseñar la documentación metodológica, funcional y de modelado del \ac{SGDIC},
que integre la gestión de cupos, materias, comunicación, faltas, evaluación
docente, quejas y sugerencias, junto con un componente de analítica de datos que
genere soluciones rápidas y verificables para el \ac{DIC}.
\end{cajaObjetivo}

\subsection{Objetivos específicos}
\begin{enumerate}[leftmargin=1.4cm, itemsep=4pt]
  \item Documentar el problema y los procesos actuales del \ac{DIC} mediante
        técnicas de investigación cualitativa y cuantitativa.
  \item Definir la misión, visión, valores y objetivos estratégicos del sistema.
  \item Especificar los requerimientos funcionales y no funcionales por actor.
  \item Modelar el sistema mediante diagramas \ac{UML}: casos de uso, clases,
        secuencia, actividades y estados.
  \item Diseñar el mecanismo de evaluación docente con control de plazos que
        garantice su aplicación a tiempo.
  \item Diseñar el módulo de cuestionarios de materias y horarios, y el modelo de
        preferencias por probabilidad.
  \item Proponer el algoritmo de soluciones rápidas basado en analítica de datos.
  \item Definir el modelo de verificación de resultados para la secretaría y el
        departamento.
  \item Establecer el cronograma, los riesgos y los indicadores de éxito.
\end{enumerate}

\section{Alcance y Delimitaciones}

\subsection{Alcance}
El proyecto cubre el análisis, la metodología de investigación, la especificación
de requerimientos y el modelado completo del sistema, así como la definición del
algoritmo analítico y del modelo de verificación institucional. Incluye:

\begin{itemize}[leftmargin=1.4cm, itemsep=3pt]
  \item Los cuatro actores: estudiante, docente, secretaría y departamento.
  \item Los once procesos problemáticos identificados en la sección anterior.
  \item Los diagramas \ac{UML} de casos de uso, clases, secuencia, actividades y
        estados, embebidos en este documento.
  \item La especificación del algoritmo de soluciones rápidas y las métricas.
\end{itemize}

\subsection{Delimitaciones}
\begin{itemize}[leftmargin=1.4cm, itemsep=3pt]
  \item No se incluye la implementación completa del sistema productivo, sino su
        documentación y modelado.
  \item El análisis se limita al \ac{DIC}; otras unidades académicas podrían
        requerir adaptaciones.
  \item Los datos usados en los ejemplos son sintéticos y anonimizados, en
        cumplimiento de las normas de protección de datos personales (\ac{PII}).
  \item La integración con el sistema central de la universidad se especifica a
        nivel de interfaz (\ac{API} \ac{REST}), sin detallar su implementación.
\end{itemize}

\section{Hipótesis}

\begin{enumerate}[leftmargin=1.4cm, itemsep=4pt]
  \item \textbf{H1.} Un canal unificado de comunicación reduce en al menos
        $60\%$ el tiempo de respuesta a solicitudes y mensajes respecto al
        proceso manual actual.
  \item \textbf{H2.} Un mecanismo de evaluación con recordatorios automáticos,
        plazos visibles e incentivos incrementa la tasa de respuesta docente
        a más del $90\%$ dentro del periodo.
  \item \textbf{H3.} Un modelo de preferencias por probabilidad sobre los
        cuestionarios de horarios permite asignar horarios con mayor nivel de
        aceptación por parte de los estudiantes.
  \item \textbf{H4.} Un algoritmo de priorización basado en analítica reduce el
        tiempo de resolución de quejas críticas y mejora la percepción de
        efectividad de la secretaría.
  \item \textbf{H5.} La trazabilidad completa (estado, responsable, fechas)
        aumenta la confianza institucional y la participación en los canales
        formales.
\end{enumerate}

\section{Metodología Resumida}

La investigación se aborda con un enfoque \textbf{mixto} (cualitativo y
cuantitativo) y un diseño de \textbf{investigación aplicada} complementado con el
ciclo \ac{CRISP-DM} para el componente de analítica de datos. El detalle completo
se presenta en el Capítulo~\ref{ch:metodologia}.

\section{Estructura del Documento}

\begin{table}[H]
\centering
\caption{Estructura de la documentación}
\label{tab:estructura}
\begin{tabularx}{\textwidth}{@{}c l X@{}}
\toprule
\rowcolor{azulClaro}
\textbf{Cap.} & \textbf{Título} & \textbf{Contenido principal} \\
\midrule
1 & Introducción y planteamiento del problema &
    Contexto, problema, preguntas, justificación, objetivos, alcance, hipótesis. \\
2 & Metodología de investigación &
    Enfoque mixto, diseño, población, instrumentos, CRISP-DM, validez y ética. \\
3 & Misión, visión y objetivos &
    Misión, visión, valores, objetivos estratégicos, \ac{KPI} y mapa estratégico. \\
4 & Requerimientos &
    Actores, requerimientos funcionales y no funcionales, reglas de negocio. \\
5 & Modelado \ac{UML} &
    Casos de uso, clases, secuencia, actividades y estados (diagramas embebidos). \\
6 & Algoritmo de soluciones rápidas &
    Analítica de datos, priorización, predicción de horarios, verificación. \\
7 & Planificación &
    Cronograma, recursos, riesgos, indicadores y conclusiones. \\
\bottomrule
\end{tabularx}
\end{table}

\section{Glosario de Términos del Dominio}

\begin{table}[H]
\centering
\caption{Glosario del dominio académico}
\label{tab:glosario}
\begin{tabularx}{\textwidth}{@{}l X@{}}
\toprule
\rowcolor{verdeClaro}
\textbf{Término} & \textbf{Definición} \\
\midrule
Cupo & Plaza disponible de un estudiante en una materia para un periodo académico. \\
Materia & Unidad curricular con código, créditos, prerrequisitos y horario. \\
Falta & Inasistencia de un estudiante a una sesión de clase programada. \\
Queja & Manifestación formal de insatisfacción sobre un servicio o persona. \\
Sugerencia & Propuesta de mejora no vinculada a una inconformidad. \\
Evaluación docente & Instrumento mediante el cual el estudiante valora el desempeño. \\
Cuestionario de horario & Instrumento que captura preferencias y disponibilidad. \\
Solución rápida & Recomendación accionable generada por el motor analítico. \\
Verificación & Contraste de resultados, quejas y evaluaciones por secretaría y dirección. \\
Trazabilidad & Conocer el estado, responsable e historial de un caso. \\
\bottomrule
\end{tabularx}
\end{table}
% =====================================================================
% CAPÍTULO 2: METODOLOGÍA DE INVESTIGACIÓN
% =====================================================================
\chapter{Metodología de Investigación}
\label{ch:metodologia}

\section{Enfoque de Investigación}

La presente investigación adopta un enfoque \textbf{mixto} (cualitativo y
cuantitativo) con alcance \textbf{descriptivo, explicativo y propositivo}. Se
eligió este enfoque porque el problema del \ac{DIC} combina dimensiones humanas
(percepciones, comunicación, conflictos) con dimensiones medibles (tiempos de
respuesta, tasas de respuesta, frecuencias de quejas, disponibilidad horaria).

\begin{table}[H]
\centering
\caption{Componentes del enfoque mixto}
\label{tab:enfoque-mixto}
\begin{tabularx}{\textwidth}{@{}l X X@{}}
\toprule
\rowcolor{azulClaro}
\textbf{Dimensión} & \textbf{Enfoque cualitativo} & \textbf{Enfoque cuantitativo} \\
\midrule
Propósito & Comprender el problema y sus causas & Medir magnitud, frecuencia e impacto \\
Técnicas & Entrevistas, observación, análisis documental & Encuestas, registro de tiempos, analítica \\
Instrumentos & Guion de entrevista, ficha de observación & Cuestionario estructurado, \ac{ETL} de logs \\
Muestra & 8--12 informantes clave & Muestreo probabilístico por estrato \\
Análisis & Codificación temática & Estadística descriptiva e inferencial \\
Resultado & Marco interpretativo del problema & Indicadores y modelos predictivos \\
\bottomrule
\end{tabularx}
\end{table}

\section{Tipo y Diseño de Investigación}

\subsection{Tipo}
\begin{itemize}[leftmargin=1.4cm, itemsep=3pt]
  \item \textbf{Investigación aplicada:} busca resolver un problema concreto del
        \ac{DIC} mediante una solución tecnológica documentada.
  \item \textbf{Estudio de caso:} el \ac{DIC} se analiza como unidad integral,
        considerando sus procesos, actores y datos.
  \item \textbf{Investigación--acción:} se iteran ciclos de diagnóstico, diseño,
        validación y refinamiento con participación de los actores.
\end{itemize}

\subsection{Diseño}
Se emplea un diseño \textbf{transversal} para el diagnóstico (una medición
completa del estado actual) y \textbf{longitudinal} para la validación (mediciones
del antes y después de la implementación de los módulos).

\begin{figure}[H]
\centering
\begin{tikzpicture}[node distance=0.55cm and 1.1cm, every node/.style={font=\scriptsize}]

\node[terminal] (inicio) {Inicio};
\node[flujo, below=of inicio] (f1) {Diagnóstico del\\estado actual};
\node[flujo, below=of f1] (f2) {Recolección cualitativa\\y cuantitativa};
\node[flujo, below=of f2] (f3) {Análisis y\\codificación};
\node[decision, below=of f3, yshift=-0.35cm] (d1) {¿Problema\\validado?};
\node[flujo, left=of d1, xshift=-0.9cm] (f4) {Ampliar\\muestra};
\node[flujo, below=of d1, yshift=-0.5cm] (f5) {Diseño de\\la solución};
\node[flujo, below=of f5] (f6) {Prototipo y\\validación};
\node[decision, below=of f6, yshift=-0.35cm] (d2) {¿Hipótesis\\aceptada?};
\node[flujo, right=of d2, xshift=0.9cm] (f7) {Refinar\\iterativamente};
\node[terminal, below=of d2, yshift=-0.5cm] (fin) {Documentación\\final};

\draw[flecha] (inicio) -- (f1);
\draw[flecha] (f1) -- (f2);
\draw[flecha] (f2) -- (f3);
\draw[flecha] (f3) -- (d1);
\draw[flecha] (d1) -- node[above, font=\tiny]{No} (f4);
\draw[flecha] (f4) |- (f2);
\draw[flecha] (d1) -- node[right, font=\tiny]{Sí} (f5);
\draw[flecha] (f5) -- (f6);
\draw[flecha] (f6) -- (d2);
\draw[flecha] (d2) -- node[above, font=\tiny]{No} (f7);
\draw[flecha] (f7) |- (f5);
\draw[flecha] (d2) -- node[right, font=\tiny]{Sí} (fin);

\end{tikzpicture}
\caption{Ciclo metodológico de investigación--acción aplicado al \ac{DIC}}
\label{fig:ciclo-metodologico}
\end{figure}

\section{Fases de la Metodología}

La metodología se estructura en ocho fases, cada una con entradas, actividades y
entregables explícitos.

\begin{table}[H]
\centering
\caption{Fases de la metodología de investigación}
\label{tab:fases}
\small
\begin{tabularx}{\textwidth}{@{}c l X l@{}}
\toprule
\rowcolor{azulClaro}
\textbf{Fase} & \textbf{Nombre} & \textbf{Actividades principales} & \textbf{Duración} \\
\midrule
F1 & Definición del problema &
    Identificación de síntomas, delimitación del alcance, formulación de preguntas
    e hipótesis. & 2 sem. \\
F2 & Revisión de literatura &
    Búsqueda de antecedentes en gestión académica, analítica educativa, sistemas
    de quejas y evaluación docente. & 3 sem. \\
F3 & Diseño metodológico &
    Selección de técnicas, definición de población y muestra, construcción de
    instrumentos. & 2 sem. \\
F4 & Recolección de datos &
    Entrevistas, encuestas, observación de procesos, extracción de datos
    históricos del \ac{DIC}. & 4 sem. \\
F5 & Análisis de datos &
    Transcripción y codificación, estadística descriptiva, análisis de
    correlaciones, preparación del dataset. & 3 sem. \\
F6 & Diseño de la solución &
    Arquitectura, modelado \ac{UML}, especificación de requerimientos y del
    algoritmo de soluciones rápidas. & 4 sem. \\
F7 & Validación &
    Validación de instrumentos, validación del modelo con expertos, prueba piloto
    del prototipo. & 3 sem. \\
F8 & Documentación y difusión &
    Redacción del informe, tableros de resultados, socialización con la comunidad
    académica. & 3 sem. \\
\bottomrule
\end{tabularx}
\end{table}

\section{Población y Muestra}

\subsection{Población}
La población objeto de estudio está conformada por los integrantes del \ac{DIC}:

\begin{itemize}[leftmargin=1.4cm, itemsep=3pt]
  \item \textbf{Estudiantes:} matriculados en materias del departamento en el
        periodo académico vigente.
  \item \textbf{Docentes:} adscritos al departamento.
  \item \textbf{Personal de secretaría:} responsable de la atención de
        solicitudes, quejas y sugerencias.
  \item \textbf{Directivos:} dirección y coordinación académica del departamento.
\end{itemize}

\subsection{Tamaño de muestra}
Para la encuesta a estudiantes se utiliza muestreo probabilístico estratificado
por semestre, con un nivel de confianza del $95\%$ y un margen de error del $5\%$:

\begin{equation}
n = \frac{N \cdot Z^{2} \cdot p \cdot q}{e^{2}(N-1) + Z^{2} \cdot p \cdot q}
\label{eq:muestra}
\end{equation}

donde $N$ es el tamaño de la población, $Z = 1.96$, $p = q = 0.5$ y
$e = 0.05$. Para el resto de actores se emplea \textbf{muestreo no probabilístico
por conveniencia}, dado el tamaño reducido de la población (docentes, secretaría
y directivos).

\begin{table}[H]
\centering
\caption{Distribución de la muestra por actor}
\label{tab:muestra}
\begin{tabularx}{\textwidth}{@{}l c c X@{}}
\toprule
\rowcolor{verdeClaro}
\textbf{Actor} & \textbf{Población} & \textbf{Muestra} & \textbf{Técnica principal} \\
\midrule
Estudiantes & $N$ (variable) & $n$ calculado & Encuesta estructurada \\
Docentes & 20--40 & 12--20 & Entrevista semiestructurada \\
Secretaría & 2--4 & 2--4 & Entrevista + observación de procesos \\
Directivos & 2--3 & 2--3 & Entrevista en profundidad \\
Datos históricos & Registros del \ac{DIC} & Censo & Extracción y \ac{ETL} \\
\bottomrule
\end{tabularx}
\end{table}

\section{Técnicas e Instrumentos de Recolección}

\begin{table}[H]
\centering
\caption{Técnicas e instrumentos de recolección de datos}
\label{tab:instrumentos}
\small
\begin{tabularx}{\textwidth}{@{}p{3.4cm} p{4.9cm} X@{}}
\toprule
\rowcolor{naranjaClaro}
\textbf{Técnica} & \textbf{Instrumento} & \textbf{Propósito} \\
\midrule
Encuesta & Cuestionario de 25 ítems (escala Likert 1--5) &
    Medir percepción de tiempos, satisfacción con comunicación y utilidad percibida \\
Entrevista semiestructurada & Guion de 15 preguntas &
    Comprender causas del retraso, sobrecarga operativa y expectativas \\
Observación no participante & Ficha de registro de procesos &
    Cronometrar el flujo real de una solicitud de cupo y de una queja \\
Análisis documental & Ficha de revisión normativa &
    Identificar reglas de negocio y criterios de asignación vigentes \\
Cuestionario de horarios & Formulario de disponibilidad y preferencia &
    Capturar datos para el modelo probabilístico de horarios \\
Extracción de datos & Consultas sobre registros históricos &
    Construir el dataset para analítica de datos \\
Grupo focal & Guion de discusión (6--8 estudiantes) &
    Validar hallazgos y priorizar funcionalidades \\
\bottomrule
\end{tabularx}
\end{table}

\subsection{Escala de medición}
Se emplea una escala Likert de cinco niveles para las variables perceptuales:

\begin{table}[H]
\centering
\caption{Escala Likert utilizada en los cuestionarios}
\label{tab:likert}
\begin{tabular}{@{}c l@{}}
\toprule
\rowcolor{moradoClaro}
\textbf{Valor} & \textbf{Significado} \\
\midrule
5 & Totalmente de acuerdo / Muy satisfecho \\
4 & De acuerdo / Satisfecho \\
3 & Neutral / Indiferente \\
2 & En desacuerdo / Insatisfecho \\
1 & Totalmente en desacuerdo / Muy insatisfecho \\
\bottomrule
\end{tabular}
\end{table}

\subsection{Operacionalización de variables}
\begin{table}[H]
\centering
\caption{Operacionalización de variables}
\label{tab:variables}
\small
\begin{tabularx}{\textwidth}{@{}l l l X@{}}
\toprule
\rowcolor{azulClaro}
\textbf{Variable} & \textbf{Tipo} & \textbf{Indicador} & \textbf{Instrumento} \\
\midrule
Tiempo de respuesta & Cuantitativa continua & Horas promedio & Observación, logs \\
Tasa de respuesta docente & Cuantitativa discreta & Porcentaje & Registro del sistema \\
Satisfacción comunicacional & Cualitativa ordinal & Índice 1--5 & Encuesta \\
Frecuencia de quejas & Cuantitativa discreta & N.º por periodo & Registro del sistema \\
Preferencia de horario & Cualitativa nominal & Probabilidad $P(h_i)$ & Cuestionario de horarios \\
Desempeño docente & Cuantitativa continua & Puntaje 0--5 & Evaluación docente \\
Efectividad secretaría & Cuantitativa continua & $\%$ casos resueltos & Tablero verificación \\
\bottomrule
\end{tabularx}
\end{table}

\section{Metodología CRISP-DM para la Analítica de Datos}

Para el componente de analítica de datos y el algoritmo de soluciones rápidas se
adopta el estándar \ac{CRISP-DM}, que organiza el trabajo de datos en seis fases
iterativas.

\begin{figure}[H]
\centering
\begin{tikzpicture}[
  fase/.style={
    draw=azulOscuro, fill=azulClaro, thick, rounded corners=6pt,
    text=azulOscuro, align=center, font=\small\bfseries,
    minimum width=3.2cm, minimum height=1.1cm, inner sep=4pt
  },
  node distance=0.7cm and 0.9cm
]

\node[fase] (f1) {1. Comprensión\\del negocio};
\node[fase, right=of f1] (f2) {2. Comprensión\\de los datos};
\node[fase, right=of f2] (f3) {3. Preparación\\de los datos};
\node[fase, below=of f3] (f4) {4. Modelado};
\node[fase, left=of f4] (f5) {5. Evaluación};
\node[fase, left=of f5] (f6) {6. Despliegue};

\draw[flecha] (f1) -- (f2);
\draw[flecha] (f2) -- (f3);
\draw[flecha] (f3) -- (f4);
\draw[flecha] (f4) -- (f5);
\draw[flecha] (f5) -- (f6);
\draw[flecha] (f6) -- (f1);
\draw[-{Stealth[length=2mm,width=1.4mm]}, thick, dashed, color=moradoUC, bend right=35]
      (f4) to node[below, font=\tiny]{iterar} (f3);
\draw[-{Stealth[length=2mm,width=1.4mm]}, thick, dashed, color=moradoUC, bend left=35]
      (f5) to node[above, font=\tiny]{refinar} (f2);

\end{tikzpicture}
\caption{Ciclo \ac{CRISP-DM} adaptado a la analítica del \ac{SGDIC}}
\label{fig:crispdm}
\end{figure}

\begin{table}[H]
\centering
\caption{Aplicación de \ac{CRISP-DM} al \ac{SGDIC}}
\label{tab:crispdm}
\small
\begin{tabularx}{\textwidth}{@{}c l X@{}}
\toprule
\rowcolor{verdeClaro}
\textbf{Fase} & \textbf{Nombre} & \textbf{Aplicación concreta en el \ac{DIC}} \\
\midrule
1 & Comprensión del negocio &
    Objetivos: reducir tiempos de respuesta, asegurar evaluación docente
    oportuna, priorizar quejas. Definición de \ac{KPI}. \\
2 & Comprensión de los datos &
    Fuentes: registros de cupos, mensajes, faltas, evaluaciones, quejas,
    cuestionarios de horarios. \\
3 & Preparación de los datos &
    Limpieza, normalización de fechas y categorías, anonimización de \ac{PII},
    construcción de variables derivadas. \\
4 & Modelado &
    Modelos de clasificación de quejas, predicción de demanda de cupos,
    probabilidad de preferencia de horario, detección de anomalías. \\
5 & Evaluación &
    Validación cruzada, matrices de confusión, comparación contra línea base,
    revisión con expertos del departamento. \\
6 & Despliegue &
    Integración de las recomendaciones en el tablero de verificación y en el
    motor de soluciones rápidas. \\
\bottomrule
\end{tabularx}
\end{table}

\section{Validez y Confiabilidad}

\subsection{Validez de contenido}
Los instrumentos se someten a \textbf{juicio de expertos} (tres docentes con
experiencia en metodología de investigación y dos directivos del \ac{DIC}), con
una rúbrica de evaluación de claridad, pertinencia y suficiencia.

\subsection{Validez de constructo}
Se aplica \textbf{análisis factorial exploratorio} para verificar que los ítems
del cuestionario agrupan las dimensiones previstas: comunicación, oportunidad,
trazabilidad y satisfacción.

\subsection{Confiabilidad}
La consistencia interna se mide con el coeficiente \textbf{alfa de Cronbach}:

\begin{equation}
\alpha = \frac{k}{k-1}\left(1 - \frac{\sum_{i=1}^{k} \sigma^{2}_{i}}{\sigma^{2}_{T}}\right)
\label{eq:cronbach}
\end{equation}

donde $k$ es el número de ítems, $\sigma^{2}_{i}$ la varianza de cada ítem y
$\sigma^{2}_{T}$ la varianza total. Se considera aceptable un valor
$\alpha \geq 0.70$.

\subsection{Validez de la analítica}
Para los modelos de datos se emplea \textbf{validación cruzada de $k$ particiones}
($k=10$) y se reportan exactitud, precisión, exhaustividad y $F_1$:

\begin{equation}
F_1 = 2 \cdot \frac{\text{Precisión} \cdot \text{Exhaustividad}}
{\text{Precisión} + \text{Exhaustividad}}
\label{eq:f1}
\end{equation}

\section{Consideraciones Éticas}

\begin{cajaNota}[Principios éticos del estudio]
\begin{itemize}[leftmargin=1.2cm, itemsep=3pt]
  \item \textbf{Consentimiento informado:} todos los participantes conocen el
        propósito del estudio y autorizan su participación.
  \item \textbf{Confidencialidad:} las respuestas de estudiantes sobre docentes
        se anonimizan; ningún docente conoce la identidad de quien lo califica.
  \item \textbf{No represalia:} se garantiza que reportar una queja o una
        evaluación negativa no genere consecuencias académicas adversas.
  \item \textbf{Minimización de datos:} solo se recolectan los datos necesarios
        para cada objetivo declarado.
  \item \textbf{Transparencia algorítmica:} los criterios del algoritmo de
        priorización son públicos y auditables por la comunidad académica.
  \item \textbf{Revisión institucional:} el protocolo se somete al comité de
        ética o instancia equivalente de la facultad.
\end{itemize}
\end{cajaNota}

\section{Análisis de Datos Previsto}

\begin{table}[H]
\centering
\caption{Técnicas de análisis por tipo de dato}
\label{tab:analisis}
\small
\begin{tabularx}{\textwidth}{@{}p{3.4cm} p{5.4cm} X@{}}
\toprule
\rowcolor{azulClaro}
\textbf{Tipo de dato} & \textbf{Técnica} & \textbf{Resultado esperado} \\
\midrule
Respuestas Likert & Estadística descriptiva (media, mediana, desviación) &
    Índice de satisfacción por dimensión \\
Comparación antes/después & Prueba $t$ para muestras relacionadas &
    Significancia del cambio en tiempos de respuesta \\
Datos categóricos de quejas & Tablas de contingencia, $\chi^{2}$ &
    Asociación entre tipo de queja y área responsable \\
Series temporales de cupos & Suavizado exponencial, ARIMA &
    Pronóstico de demanda por materia \\
Preferencias de horario & Frecuencias y probabilidad condicional &
    $P(\text{horario} \mid \text{perfil del estudiante})$ \\
Texto libre de quejas & PLN (bolsa de palabras) &
    Categorización automática y temas recurrentes \\
Registros de evaluación & Tendencia y control estadístico &
    Detección de docentes con desempeño atípico \\
\bottomrule
\end{tabularx}
\end{table}

\section{Limitaciones Metodológicas}

\begin{itemize}[leftmargin=1.4cm, itemsep=3pt]
  \item \textbf{Sesgo de autoselección:} quienes responden encuestas pueden tener
        opiniones más extremas que quienes no participan.
  \item \textbf{Deseabilidad social:} los estudiantes podrían suavizar sus
        evaluaciones por temor a represalias, mitigado con anonimato estricto.
  \item \textbf{Datos históricos incompletos:} los registros previos pueden ser
        fragmentarios, lo que obliga a imputar o descartar observaciones.
  \item \textbf{Estacionalidad académica:} los patrones de quejas y demanda de
        cupos varían según el calendario académico, lo que exige comparar
        periodos equivalentes.
  \item \textbf{Generalización:} los hallazgos describen el \ac{DIC} y no se
        extrapolan automáticamente a otras unidades académicas.
\end{itemize}

% =====================================================================
% CAPÍTULO 3: MISIÓN, VISIÓN, VALORES Y OBJETIVOS
% =====================================================================
\chapter{Misión, Visión, Valores y Objetivos}
\label{ch:mision-vision}

\section{Identidad del Sistema}

\begin{cajaNota}[Identidad]
\textbf{Denominación:} Sistema de Gestión Departamental de Informática y
Computación (\sistema).\\[4pt]
\textbf{Naturaleza:} Plataforma integrada de gestión académica, comunicación
institucional y analítica de datos.\\[4pt]
\textbf{Ámbito de aplicación:} \ac{DIC}.\\[4pt]
\textbf{Usuarios:} estudiantes, docentes, secretaría y dirección del departamento.
\end{cajaNota}

\section{Misión}

\begin{cajaObjetivo}[Misión]
Facilitar al \ac{DIC} una plataforma unificada, transparente y trazable que
centralice la gestión de cupos, materias, comunicación, faltas, evaluación
docente, quejas y sugerencias, incorporando analítica de datos que permita
detectar patrones, anticipar problemas y generar soluciones rápidas basadas en
evidencia, con el fin de mejorar la oportunidad de las respuestas, la equidad en
las decisiones y la calidad de la experiencia académica de estudiantes, docentes
y personal administrativo.
\end{cajaObjetivo}

\subsection{Desglose de la misión}
\begin{table}[H]
\centering
\caption{Componentes de la misión}
\label{tab:mision-desglose}
\begin{tabularx}{\textwidth}{@{}l X@{}}
\toprule
\rowcolor{verdeClaro}
\textbf{Componente} & \textbf{Contenido} \\
\midrule
¿Qué hacemos? & Centralizamos y automatizamos los procesos académicos y
administrativos del \ac{DIC}. \\
¿Para quién? & Estudiantes, docentes, secretaría y dirección del departamento. \\
¿Para qué? & Para mejorar oportunidad, trazabilidad, equidad y calidad de las
decisiones académicas. \\
¿Cómo? & Mediante una plataforma integrada con analítica de datos y un motor de
soluciones rápidas. \\
¿Con qué valores? & Transparencia, equidad, confidencialidad, eficiencia,
innovación y servicio. \\
\bottomrule
\end{tabularx}
\end{table}

\section{Visión}

\begin{cajaNota}[Visión]
Ser, en un horizonte de tres años, el sistema de gestión departamental de
referencia en la universidad, reconocido por eliminar las barreras de
comunicación entre estudiantes y docentes, por garantizar evaluaciones docentes
oportunas y por transformar los datos operativos en decisiones académicas
informadas, anticipadas y verificables.
\end{cajaNota}

\subsection{Metas de la visión}
\begin{table}[H]
\centering
\caption{Metas cuantificables asociadas a la visión}
\label{tab:metas-vision}
\small
\begin{tabularx}{\textwidth}{@{}X c c@{}}
\toprule
\rowcolor{azulClaro}
\textbf{Meta} & \textbf{Línea base} & \textbf{Objetivo a 3 años} \\
\midrule
Tiempo de respuesta a solicitudes de cupo & 5--10 días hábiles & $< 24$ horas \\
Evaluaciones docentes completadas en el periodo & $< 40\%$ & $= 100\%$ \\
Quejas con responsable asignado y fecha de cierre & $< 30\%$ & $> 95\%$ \\
Tiempo medio de resolución de queja crítica & $> 15$ días & $< 3$ días \\
Cobertura de la encuesta de horarios & No existente & $> 80\%$ de matriculados \\
Módulos con tablero de verificación & 0 & 100\% de procesos críticos \\
Materias con horario optimizado por modelo & 0 & $> 70\%$ \\
\bottomrule
\end{tabularx}
\end{table}

\section{Valores Institucionales}

\begin{table}[H]
\centering
\caption{Valores y su manifestación en el sistema}
\label{tab:valores}
\small
\begin{tabularx}{\textwidth}{@{}l X@{}}
\toprule
\rowcolor{moradoClaro}
\textbf{Valor} & \textbf{Cómo se manifiesta en el \ac{SGDIC}} \\
\midrule
Transparencia & Estados, plazos y criterios visibles para el solicitante. \\
Equidad & Reglas de priorización publicadas y aplicadas uniformemente. \\
Confidencialidad & Anonimato del evaluador y protección de datos personales. \\
Trazabilidad & Cada caso registra responsable, fecha y bitácora de cambios. \\
Oportunidad & Notificaciones automáticas y control de plazos. \\
Eficiencia & Automatización y eliminación de intermediarios innecesarios. \\
Innovación & Uso de analítica y modelos predictivos para anticipar problemas. \\
Servicio & Foco en reducir la carga y mejorar la experiencia de los actores. \\
Mejora continua & Ciclos de retroalimentación a partir de las métricas. \\
\bottomrule
\end{tabularx}
\end{table}

\section{Objetivos Estratégicos}

\begin{enumerate}[leftmargin=1.4cm, itemsep=5pt]
  \item \textbf{OE1. Gestión de cupos eficiente.} Automatizar la solicitud,
        verificación de disponibilidad y asignación de cupos con criterios
        transparentes y tiempo de respuesta inferior a 24 horas.

  \item \textbf{OE2. Ciclo de vida de materias.} Formalizar el flujo de creación,
        edición, aprobación y publicación de materias con control de versiones.

  \item \textbf{OE3. Comunicación institucional unificada.} Proveer un canal
        único y trazable para la comunicación bidireccional entre estudiantes,
        docentes y secretaría.

  \item \textbf{OE4. Notificación temprana de faltas.} Registrar faltas en el
        momento y notificar al estudiante de forma automática e inmediata.

  \item \textbf{OE5. Evaluación docente oportuna.} Garantizar la aplicación,
        recordatorio y cierre de la evaluación docente dentro del periodo.

  \item \textbf{OE6. Gestión formal de quejas y sugerencias.} Categorizar,
        priorizar, asignar responsable y cerrar cada caso con trazabilidad total.

  \item \textbf{OE7. Fortalecimiento de la secretaría.} Dotar a la secretaría de
        herramientas de priorización, delegación y comunicación con los docentes.

  \item \textbf{OE8. Optimización de horarios.} Aplicar cuestionarios y modelos
        probabilísticos para programar horarios con mayor aceptación estudiantil.

  \item \textbf{OE9. Verificación institucional.} Ofrecer tableros de control a
        la secretaría y al departamento para verificar resultados, quejas y
        evaluaciones.

  \item \textbf{OE10. Analítica y soluciones rápidas.} Implementar un motor
        analítico que detecte anomalías y genere recomendaciones accionables.
\end{enumerate}

\section{Indicadores Clave de Desempeño}

\begin{table}[H]
\centering
\caption{Indicadores clave de desempeño (\ac{KPI})}
\label{tab:kpis}
\small
\begin{tabularx}{\textwidth}{@{}p{1.5cm} X p{1.7cm} p{2.0cm}@{}}
\toprule
\rowcolor{azulClaro}
\textbf{KPI} & \textbf{Definición} & \textbf{Meta} & \textbf{Frecuencia} \\
\midrule
TRC & Tiempo de respuesta a solicitud de cupo (horas) & $<24$ & Semanal \\
TRD & Tiempo de respuesta docente a mensaje (horas) & $<48$ & Semanal \\
TED & Tasa de evaluaciones docentes completadas & $100\%$ & Por periodo \\
TQR & Tasa de quejas resueltas dentro del plazo & $>90\%$ & Mensual \\
TMC & Tiempo medio de cierre de queja crítica (días) & $<3$ & Mensual \\
ICOM & Índice de satisfacción con la comunicación (1--5) & $>4.0$ & Por periodo \\
CNF & Cobertura de notificación de faltas el mismo día & $100\%$ & Semanal \\
ACE & Aceptación de horario asignado (1--5) & $>4.0$ & Por periodo \\
IDD & Índice de desempeño docente promedio (0--5) & $>4.0$ & Por periodo \\
EFA & Efectividad de acciones de la secretaría & $>85\%$ & Mensual \\
EAN & Exactitud del modelo de detección de anomalías & $>85\%$ & Trimestral \\
\bottomrule
\end{tabularx}
\end{table}

\section{Mapa Estratégico}

\begin{figure}[H]
\centering
\begin{tikzpicture}[
  cajita/.style={
    draw=azulOscuro, fill=azulClaro, thick, rounded corners=4pt,
    text=azulOscuro, align=center, font=\scriptsize,
    minimum width=7.5cm, minimum height=1.2cm, inner sep=4pt
  },
  meta/.style={
    draw=verdeUC, fill=verdeClaro, thick, rounded corners=4pt,
    text=verdeUC!60!black, align=center, font=\scriptsize,
    minimum width=7.5cm, minimum height=1.2cm, inner sep=4pt
  },
  node distance=0.5cm
]

\node[meta] (vision) {\textbf{VISIÓN}\\Ser referente universitario en gestión departamental};
\node[cajita, below=of vision] (mision) {\textbf{MISIÓN}\\Plataforma unificada, trazable y analítica};
\node[cajita, below=of mision] (perspectiva) {\textbf{PERSPECTIVA ESTRATÉGICA}\\Soluciones rápidas basadas en evidencia};
\node[cajita, below=of perspectiva] (resultados) {\textbf{RESULTADOS}\\KPI: TRC, TED, TQR, CNF, ACE, IDD};

\draw[flecha] (vision) -- (mision);
\draw[flecha] (mision) -- (perspectiva);
\draw[flecha] (perspectiva) -- (resultados);

\end{tikzpicture}
\caption{Mapa estratégico del \ac{SGDIC}: de la visión a los resultados medibles}
\label{fig:mapa-estrategico}
\end{figure}

\section{Beneficios Esperados por Actor}

\begin{table}[H]
\centering
\caption{Beneficios esperados por actor involucrado}
\label{tab:beneficios}
\small
\begin{tabularx}{\textwidth}{@{}l X@{}}
\toprule
\rowcolor{verdeClaro}
\textbf{Actor} & \textbf{Beneficios principales} \\
\midrule
Estudiante &
  Solicitud de cupos en línea con respuesta en menos de 24 h; notificación
  inmediata de faltas; canal formal y anónimo para evaluar y reportar; horarios
  más compatibles con su disponibilidad. \\
Docente &
  Gestión ágil de materias; canal único de mensajería; retroalimentación de
  estudiantes consolidada y anónima; menor carga administrativa. \\
Secretaría &
  Bandeja priorizada de solicitudes y quejas; asignación de responsables;
  comunicación formal con docentes; tableros de seguimiento; eliminación de
  registros duplicados en hojas de cálculo. \\
Departamento &
  Verificación consolidada de resultados, quejas y evaluaciones; indicadores en
  tiempo real; alertas tempranas; soporte objetivo para decisiones académicas. \\
Institución &
  Mayor satisfacción de la comunidad académica; trazabilidad auditable;
  experiencia replicable en otras unidades académicas. \\
\bottomrule
\end{tabularx}
\end{table}
% =====================================================================
% CAPÍTULO 4: REQUERIMIENTOS
% =====================================================================
\chapter{Requerimientos del Sistema}
\label{ch:requerimientos}

\section{Actores del Sistema}

\begin{table}[H]
\centering
\caption{Definición de actores del \ac{SGDIC}}
\label{tab:actores}
\small
\begin{tabularx}{\textwidth}{@{}l X@{}}
\toprule
\rowcolor{azulClaro}
\textbf{Actor} & \textbf{Descripción y responsabilidades} \\
\midrule
Estudiante &
  Usuario final que solicita cupos, consulta horarios, envía y recibe mensajes,
  responde cuestionarios de horario y de materias, califica docentes, recibe
  notificaciones de faltas y presenta quejas o sugerencias. \\
Docente &
  Gestiona sus materias y horarios, registra faltas, se comunica con estudiantes
  y con la secretaría, recibe evaluaciones consolidadas. \\
Secretaría &
  Recibe y clasifica quejas y solicitudes, asigna responsables, coordina la
  comunicación con docentes, verifica resultados y mantiene la trazabilidad. \\
Departamento &
  Administra el catálogo de materias y cupos, configura periodos de evaluación,
  consulta analítica, verifica resultados y aprueba decisiones académicas. \\
Sistema de notificaciones &
  Actor secundario que despacha correos, mensajes internos y recordatorios. \\
Sistema central universitario &
  Actor secundario que provee la matrícula, catálogo oficial y autenticación. \\
\bottomrule
\end{tabularx}
\end{table}

\section{Requerimientos Funcionales}

\subsection{Módulo de autenticación y roles}

\begin{longtable}{@{}p{1.5cm} p{6.5cm} p{4.5cm}@{}}
\caption{Requerimientos funcionales --- autenticación y roles}
\label{tab:rf-auth}\\
\toprule
\rowcolor{azulClaro}
\textbf{ID} & \textbf{Requerimiento} & \textbf{Actor} \\
\midrule
\endfirsthead
\toprule
\rowcolor{azulClaro}
\textbf{ID} & \textbf{Requerimiento} & \textbf{Actor} \\
\midrule
\endhead
RF-01 & El sistema debe autenticar usuarios mediante correo institucional y contraseña o credencial institucional. & Todos \\
RF-02 & El sistema debe asignar permisos según el rol (\ac{RBAC}): estudiante, docente, secretaría, departamento. & Todos \\
RF-03 & El sistema debe permitir la recuperación de contraseña sin intervención de la secretaría. & Todos \\
RF-04 & El sistema debe registrar bitácora de accesos y operaciones críticas. & Departamento \\
\bottomrule
\end{longtable}

\subsection{Módulo de cupos}

\begin{longtable}{@{}p{1.5cm} p{6.5cm} p{4.5cm}@{}}
\caption{Requerimientos funcionales --- gestión de cupos}
\label{tab:rf-cupos}\\
\toprule
\rowcolor{verdeClaro}
\textbf{ID} & \textbf{Requerimiento} & \textbf{Actor} \\
\midrule
\endfirsthead
\toprule
\rowcolor{verdeClaro}
\textbf{ID} & \textbf{Requerimiento} & \textbf{Actor} \\
\midrule
\endhead
RF-05 & El estudiante debe poder solicitar cupo en una materia mostrando la disponibilidad en tiempo real. & Estudiante \\
RF-06 & El sistema debe validar prerrequisitos, créditos y conflictos de horario antes de aceptar la solicitud. & Sistema \\
RF-07 & El sistema debe aplicar una regla de priorización configurable (promedio, avance curricular, orden de solicitud). & Departamento \\
RF-08 & El sistema debe notificar al estudiante el resultado de la solicitud (aprobada, en lista de espera, rechazada). & Sistema \\
RF-09 & El sistema debe liberar automáticamente cupos de solicitudes no confirmadas en el plazo definido. & Sistema \\
RF-10 & La secretaría debe poder ver y gestionar la lista de espera por materia. & Secretaría \\
\bottomrule
\end{longtable}

\subsection{Módulo de materias}

\begin{longtable}{@{}p{1.5cm} p{6.5cm} p{4.5cm}@{}}
\caption{Requerimientos funcionales --- materias}
\label{tab:rf-materias}\\
\toprule
\rowcolor{naranjaClaro}
\textbf{ID} & \textbf{Requerimiento} & \textbf{Actor} \\
\midrule
\endfirsthead
\toprule
\rowcolor{naranjaClaro}
\textbf{ID} & \textbf{Requerimiento} & \textbf{Actor} \\
\midrule
\endhead
RF-11 & El docente debe poder proponer la creación de una materia nueva con su justificación. & Docente \\
RF-12 & El docente debe poder solicitar la edición de una materia existente (contenido, créditos, prerrequisitos, horario). & Docente \\
RF-13 & El sistema debe gestionar el flujo de aprobación: propuesta, revisión del departamento, aprobación o rechazo, publicación. & Departamento \\
RF-14 & El sistema debe mantener historial de versiones de cada materia con autor y fecha del cambio. & Sistema \\
RF-15 & El sistema debe publicar automáticamente las materias aprobadas en el catálogo visible para los estudiantes. & Sistema \\
\bottomrule
\end{longtable}

\subsection{Módulo de mensajería}

\begin{longtable}{@{}p{1.5cm} p{6.5cm} p{4.5cm}@{}}
\caption{Requerimientos funcionales --- mensajería}
\label{tab:rf-mensajes}\\
\toprule
\rowcolor{moradoClaro}
\textbf{ID} & \textbf{Requerimiento} & \textbf{Actor} \\
\midrule
\endfirsthead
\toprule
\rowcolor{moradoClaro}
\textbf{ID} & \textbf{Requerimiento} & \textbf{Actor} \\
\midrule
\endhead
RF-16 & El estudiante debe poder enviar mensajes a docentes de las materias en las que está inscrito. & Estudiante \\
RF-17 & El docente debe poder enviar mensajes masivos o individuales a los estudiantes de sus materias. & Docente \\
RF-18 & El sistema debe registrar estado del mensaje: enviado, entregado y leído. & Sistema \\
RF-19 & El sistema debe alertar cuando un mensaje no tenga respuesta en el plazo configurado. & Sistema \\
RF-20 & La secretaría debe poder intervenir en un hilo cuando el caso escale. & Secretaría \\
\bottomrule
\end{longtable}

\subsection{Módulo de faltas}

\begin{longtable}{@{}p{1.5cm} p{6.5cm} p{4.5cm}@{}}
\caption{Requerimientos funcionales --- registro de faltas}
\label{tab:rf-faltas}\\
\toprule
\rowcolor{azulClaro}
\textbf{ID} & \textbf{Requerimiento} & \textbf{Actor} \\
\midrule
\endfirsthead
\toprule
\rowcolor{azulClaro}
\textbf{ID} & \textbf{Requerimiento} & \textbf{Actor} \\
\midrule
\endhead
RF-21 & El docente debe poder registrar la asistencia por sesión de clase de forma rápida (lista o código). & Docente \\
RF-22 & El sistema debe notificar al estudiante el mismo día en que se registra su falta. & Sistema \\
RF-23 & El estudiante debe poder presentar una justificación adjuntando evidencia. & Estudiante \\
RF-24 & El docente debe poder aprobar o rechazar la justificación con comentario. & Docente \\
RF-25 & El sistema debe calcular el porcentaje de inasistencia acumulado por materia. & Sistema \\
RF-26 & El sistema debe alertar al estudiante, docente y secretaría cuando el acumulado supere el umbral permitido. & Sistema \\
\bottomrule
\end{longtable}

\subsection{Módulo de evaluación docente}

\begin{longtable}{@{}p{1.5cm} p{6.5cm} p{4.5cm}@{}}
\caption{Requerimientos funcionales --- evaluación docente}
\label{tab:rf-evaluacion}\\
\toprule
\rowcolor{verdeClaro}
\textbf{ID} & \textbf{Requerimiento} & \textbf{Actor} \\
\midrule
\endfirsthead
\toprule
\rowcolor{verdeClaro}
\textbf{ID} & \textbf{Requerimiento} & \textbf{Actor} \\
\midrule
\endhead
RF-27 & El departamento debe poder configurar el periodo, las preguntas y la fecha límite de la evaluación docente. & Departamento \\
RF-28 & El sistema debe habilitar la evaluación automáticamente al inicio del periodo configurado. & Sistema \\
RF-29 & El sistema debe enviar recordatorios automáticos a los estudiantes que no hayan completado la evaluación. & Sistema \\
RF-30 & El sistema debe garantizar el anonimato total del evaluador. & Sistema \\
RF-31 & El sistema debe bloquear servicios académicos hasta completar la evaluación, si así se configura. & Sistema \\
RF-32 & El sistema debe consolidar resultados y publicarlos solo al cumplirse la fecha límite o el umbral de participación. & Sistema \\
RF-33 & El docente debe poder consultar su resultado consolidado y su comparación histórica. & Docente \\
RF-34 & El departamento debe poder identificar docentes con resultados atípicos y generar planes de acompañamiento. & Departamento \\
\bottomrule
\end{longtable}

\subsection{Módulo de quejas y sugerencias}

\begin{longtable}{@{}p{1.5cm} p{6.5cm} p{4.5cm}@{}}
\caption{Requerimientos funcionales --- quejas y sugerencias}
\label{tab:rf-quejas}\\
\toprule
\rowcolor{naranjaClaro}
\textbf{ID} & \textbf{Requerimiento} & \textbf{Actor} \\
\midrule
\endfirsthead
\toprule
\rowcolor{naranjaClaro}
\textbf{ID} & \textbf{Requerimiento} & \textbf{Actor} \\
\midrule
\endhead
RF-35 & Cualquier actor debe poder registrar una queja o sugerencia con categoría y descripción. & Todos \\
RF-36 & El sistema debe clasificar automáticamente el caso por categoría, urgencia y área responsable. & Sistema \\
RF-37 & La secretaría debe poder asignar responsable, plazo y estado a cada caso. & Secretaría \\
RF-38 & La secretaría debe poder comunicarse formalmente con el docente involucrado dentro del caso. & Secretaría \\
RF-39 & El sistema debe notificar al solicitante cada cambio de estado del caso. & Sistema \\
RF-40 & El sistema debe permitir el cierre con evidencia de la acción tomada y encuesta de satisfacción. & Secretaría \\
RF-41 & El sistema debe permitir el anonimato del denunciante cuando la naturaleza del caso lo requiera. & Sistema \\
RF-42 & El sistema debe impedir el cierre de un caso sin acción registrada. & Sistema \\
\bottomrule
\end{longtable}

\subsection{Módulo de cuestionarios de materias y horarios}

\begin{longtable}{@{}p{1.5cm} p{6.5cm} p{4.5cm}@{}}
\caption{Requerimientos funcionales --- cuestionarios}
\label{tab:rf-cuestionarios}\\
\toprule
\rowcolor{moradoClaro}
\textbf{ID} & \textbf{Requerimiento} & \textbf{Actor} \\
\midrule
\endfirsthead
\toprule
\rowcolor{moradoClaro}
\textbf{ID} & \textbf{Requerimiento} & \textbf{Actor} \\
\midrule
\endhead
RF-43 & El sistema debe aplicar un cuestionario de disponibilidad horaria por estudiante y periodo. & Estudiante \\
RF-44 & El sistema debe aplicar un cuestionario de preferencia y evaluación de materias cursadas. & Estudiante \\
RF-45 & El sistema debe calcular la probabilidad de preferencia por cada franja horaria disponible. & Sistema \\
RF-46 & El sistema debe recomendar una asignación de horario que maximice la aceptación esperada. & Departamento \\
RF-47 & El sistema debe mostrar al departamento escenarios alternativos de horario con su impacto estimado. & Departamento \\
RF-48 & El sistema debe permitir al estudiante confirmar o rechazar el horario sugerido. & Estudiante \\
\bottomrule
\end{longtable}

\subsection{Módulo de analítica y soluciones rápidas}

\begin{longtable}{@{}p{1.5cm} p{6.5cm} p{4.5cm}@{}}
\caption{Requerimientos funcionales --- analítica}
\label{tab:rf-analitica}\\
\toprule
\rowcolor{azulClaro}
\textbf{ID} & \textbf{Requerimiento} & \textbf{Actor} \\
\midrule
\endfirsthead
\toprule
\rowcolor{azulClaro}
\textbf{ID} & \textbf{Requerimiento} & \textbf{Actor} \\
\midrule
\endhead
RF-49 & El sistema debe alimentar un repositorio analítico con los datos de todos los módulos. & Sistema \\
RF-50 & El sistema debe calcular los \ac{KPI} definidos en el Capítulo~\ref{ch:mision-vision}. & Sistema \\
RF-51 & El sistema debe detectar anomalías: cupos saturados, quejas recurrentes, docentes con baja evaluación, inasistencia atípica. & Sistema \\
RF-52 & El sistema debe generar recomendaciones accionables priorizadas por impacto y esfuerzo. & Sistema \\
RF-53 & El sistema debe estimar la demanda futura de cupos por materia. & Departamento \\
RF-54 & La secretaría y el departamento deben verificar resultados, quejas y evaluaciones en tableros consolidados. & Secretaría, Departamento \\
RF-55 & El sistema debe exportar reportes a \ac{CSV} y \ac{PDF}. & Secretaría, Departamento \\
RF-56 & El sistema debe registrar si la recomendación fue aplicada y medir su efecto posterior. & Sistema \\
\bottomrule
\end{longtable}

\section{Requerimientos No Funcionales}

\begin{table}[H]
\centering
\caption{Requerimientos no funcionales}
\label{tab:rnf}
\small
\begin{tabularx}{\textwidth}{@{}p{1.5cm} l X@{}}
\toprule
\rowcolor{verdeClaro}
\textbf{ID} & \textbf{Categoría} & \textbf{Descripción} \\
\midrule
RNF-01 & Usabilidad &
    Interfaz responsiva, accesible desde móvil, con un máximo de tres clics
    para completar cualquier operación frecuente. \\
RNF-02 & Rendimiento &
    Tiempo de respuesta menor a 2 segundos en el 95\% de las peticiones. \\
RNF-03 & Disponibilidad &
    Disponibilidad mínima del 99.5\% en horario académico (06:00--22:00). \\
RNF-04 & Seguridad &
    Cifrado en tránsito (TLS) y en reposo, autenticación multifactor para roles
    administrativos. \\
RNF-05 & Privacidad &
    Anonimización de evaluaciones y quejas sensibles; cumplimiento de normas de
    protección de datos personales. \\
RNF-06 & Escalabilidad &
    Soporte para el crecimiento de usuarios y datos sin rediseño estructural. \\
RNF-07 & Mantenibilidad &
    Arquitectura modular con separación de responsabilidades e interfaces
    documentadas. \\
RNF-08 & Interoperabilidad &
    \ac{API} \ac{REST} documentada para integrarse con el sistema central
    universitario. \\
RNF-09 & Trazabilidad &
    Registro de auditoría inmutable de cada transición de estado. \\
RNF-10 & Compatibilidad &
    Funcionamiento en navegadores modernos (Chrome, Edge, Firefox, Safari). \\
RNF-11 & Accesibilidad &
    Cumplimiento de pautas WCAG 2.1 nivel AA. \\
RNF-12 & Localización &
    Interfaz en español con soporte de zona horaria y formatos locales. \\
\bottomrule
\end{tabularx}
\end{table}

\section{Reglas de Negocio}

\begin{table}[H]
\centering
\caption{Reglas de negocio del \ac{DIC}}
\label{tab:reglas}
\small
\begin{tabularx}{\textwidth}{@{}p{1.3cm} l X@{}}
\toprule
\rowcolor{naranjaClaro}
\textbf{ID} & \textbf{Nombre} & \textbf{Regla} \\
\midrule
RN-01 & Prerrequisitos &
    No se puede otorgar cupo si el estudiante no ha aprobado los prerrequisitos
    declarados en la materia. \\
RN-02 & Cupo máximo &
    El número de cupos asignados nunca puede superar el cupo máximo definido
    para la materia y el periodo. \\
RN-03 & Conflicto horario &
    No se puede asignar cupo si el horario de la materia se solapa con otra
    materia ya inscrita por el estudiante. \\
RN-04 & Priorización &
    Ante escasez de cupos, se aplica prioridad por avance curricular, luego por
    promedio acumulado y finalmente por orden cronológico de la solicitud. \\
RN-05 & Anonimato de evaluación &
    El docente nunca accede a la identidad de quien lo evaluó; solo a resultados
    agregados con un mínimo de cinco respuestas. \\
RN-06 & Ventana de evaluación &
    La evaluación docente solo puede diligenciarse dentro del periodo
    configurado; fuera de él, el formulario se bloquea. \\
RN-07 & Notificación de falta &
    Toda falta registrada debe generar notificación al estudiante el mismo día. \\
RN-08 & Umbral de inasistencia &
    Al superar el 20\% de inasistencia en una materia, se genera alerta
    automática al estudiante, al docente y a la secretaría. \\
RN-09 & Cierre de queja &
    Una queja no puede cerrarse sin registrar la acción realizada y la fecha de
    cierre. \\
RN-10 & Escalamiento &
    Una queja crítica sin atender en 3 días se escala automáticamente a la
    dirección del departamento. \\
RN-11 & Edición de materia &
    Toda modificación curricular requiere revisión y aprobación del departamento
    antes de su publicación. \\
RN-12 & Cobertura de horario &
    Un horario solo se publica si alcanza el umbral mínimo de aceptación
    estimada definido por el departamento. \\
\bottomrule
\end{tabularx}
\end{table}

\section{Matriz de Trazabilidad Objetivos--Requerimientos}

\begin{table}[H]
\centering
\caption{Trazabilidad entre objetivos estratégicos y requerimientos}
\label{tab:trazabilidad}
\small
\begin{tabularx}{\textwidth}{@{}l X@{}}
\toprule
\rowcolor{moradoClaro}
\textbf{Objetivo} & \textbf{Requerimientos asociados} \\
\midrule
OE1 (Cupos) & RF-05, RF-06, RF-07, RF-08, RF-09, RF-10, RN-01, RN-02, RN-03, RN-04 \\
OE2 (Materias) & RF-11, RF-12, RF-13, RF-14, RF-15, RN-11 \\
OE3 (Comunicación) & RF-16, RF-17, RF-18, RF-19, RF-20 \\
OE4 (Faltas) & RF-21, RF-22, RF-23, RF-24, RF-25, RF-26, RN-07, RN-08 \\
OE5 (Evaluación) & RF-27 a RF-34, RN-05, RN-06 \\
OE6 (Quejas) & RF-35 a RF-42, RN-09, RN-10 \\
OE7 (Secretaría) & RF-10, RF-20, RF-37, RF-38, RF-40, RF-54, RF-55 \\
OE8 (Horarios) & RF-43, RF-44, RF-45, RF-46, RF-47, RF-48, RN-12 \\
OE9 (Verificación) & RF-50, RF-54, RF-55, RF-56 \\
OE10 (Analítica) & RF-49, RF-51, RF-52, RF-53, RF-56 \\
\bottomrule
\end{tabularx}
\end{table}
% =====================================================================
% CAPÍTULO 5: MODELADO UML
% =====================================================================
\chapter{Modelado UML del Sistema}
\label{ch:diagramas}

El modelado se realiza con el estándar \ac{UML}~2.5. Todos los diagramas de este
capítulo están integrados directamente en el documento y se compilan junto con
el texto.

\section{Diagrama de Casos de Uso}

\subsection{Descripción general}
El diagrama de casos de uso muestra las funcionalidades del \ac{SGDIC} desde la
perspectiva de los cuatro actores principales.

\begin{figure}[H]
\centering
\resizebox{\textwidth}{!}{%
\begin{tikzpicture}[node distance=0.8cm]

% ============ FRONTERA DEL SISTEMA ============
\begin{scope}[on background layer]
\node[sistema, fit={(-2.0,0.65) (12.1,-16.2)}] (borde) {};
\end{scope}
\node[font=\small\bfseries, text=azulOscuro] at (5.05,0.25) {SGDIC --- Sistema de Gestión Departamental};

% ============ ACTORES ============
\node[actor] (est) at (-4.9,-3.0) {Estudiante};
\node[actor] (doc) at (-4.9,-9.7) {Docente};
\node[actor] (sec) at (14.3,-5.2) {Secretaría};
\node[actor] (adm) at (14.3,-11.9) {Departamento\\(Administrador)};

% ============ PAQUETE 1: GESTIÓN ACADÉMICA ============
\begin{scope}[on background layer]
\node[paquete, fit={(-1.0,-0.65) (5.0,-6.95)}] (pkg1) {};
\end{scope}
\node[font=\tiny\bfseries, text=verdeUC] at (2.0,-0.85) {Gestión Académica};
\node[caso] (uc1) at (0.5,-1.9) {Solicitar cupos};
\node[caso] (uc2) at (3.55,-1.9) {Crear/editar\\materias};
\node[caso] (uc3) at (0.5,-4.0) {Registrar faltas};
\node[caso] (uc5) at (3.55,-4.0) {Cuestionario de\\horarios y materias};
\node[caso] (uc4) at (2.0,-6.05) {Calificar docentes};

% ============ PAQUETE 2: COMUNICACIÓN ============
\begin{scope}[on background layer]
\node[paquete, fit={(-1.0,-7.35) (5.0,-12.95)}] (pkg2) {};
\end{scope}
\node[font=\tiny\bfseries, text=verdeUC] at (2.0,-7.55) {Comunicación y Retroalimentación};
\node[caso] (uc6) at (0.5,-8.65) {Mensaje de\\estudiante a docente};
\node[caso] (uc7) at (3.55,-8.65) {Mensaje de\\docente a estudiante};
\node[caso] (uc8) at (0.5,-11.05) {Quejas y\\sugerencias};
\node[caso] (uc9) at (3.55,-11.05) {Notificaciones\\automáticas};

% ============ PAQUETE 3: EVALUACIÓN Y ANÁLISIS ============
\begin{scope}[on background layer]
\node[paquete, fit={(5.8,-0.65) (11.25,-15.85)}] (pkg3) {};
\end{scope}
\node[font=\tiny\bfseries, text=verdeUC] at (8.55,-0.85) {Evaluación, Verificación y Analítica};
\node[caso] (uc10) at (8.55,-1.9) {Evaluación\\docente a tiempo};
\node[caso] (uc11) at (8.55,-4.0) {Verificar\\resultados};
\node[caso] (uc12) at (8.55,-6.1) {Ver tableros de\\analítica};
\node[caso] (uc13) at (8.55,-8.25) {Generar soluciones\\rápidas};
\node[caso] (uc14) at (8.55,-10.45) {Optimizar\\horarios};
\node[caso] (uc15) at (8.55,-12.65) {Gestionar\\actas y reportes};
\node[caso] (uc16) at (8.55,-14.8) {Administrar\\catálogo};

% ============ RELACIONES ============
\draw[thick, azulOscuro] (est) -- (uc1);
\draw[thick, azulOscuro] (est) -- (uc5);
\draw[thick, azulOscuro] (est) -- (uc4);
\draw[thick, azulOscuro] (est) -- (uc6);
\draw[thick, azulOscuro] (est) -- (uc8);

\draw[thick, azulOscuro] (doc) -- (uc2);
\draw[thick, azulOscuro] (doc) -- (uc3);
\draw[thick, azulOscuro] (doc) -- (uc7);

\draw[thick, azulOscuro] (sec) -- (uc11);
\draw[thick, azulOscuro] (sec) -- (uc12);
\draw[thick, azulOscuro] (sec) -- (uc8);
\draw[thick, azulOscuro] (sec) -- (uc15);

\draw[thick, azulOscuro] (adm) -- (uc10);
\draw[thick, azulOscuro] (adm) -- (uc13);
\draw[thick, azulOscuro] (adm) -- (uc14);
\draw[thick, azulOscuro] (adm) -- (uc16);

\draw[-{Stealth[length=2mm,width=1.4mm]}, thick, dashed, color=moradoUC]
      (uc1) -- node[above, font=\tiny, fill=white, inner sep=1pt, text=moradoUC]{$\ll$include$\gg$} (uc2);
\draw[-{Stealth[length=2mm,width=1.4mm]}, thick, dashed, color=moradoUC] (uc9) -- (uc6);
\draw[-{Stealth[length=2mm,width=1.4mm]}, thick, dashed, color=moradoUC] (uc13) -- (uc12);

\end{tikzpicture}%
}
\caption{Diagrama de casos de uso del \ac{SGDIC}}
\label{fig:casos-uso}
\end{figure}

\subsection{Especificación de casos de uso principales}

\begin{table}[H]
\centering
\caption{Especificación del caso de uso: Solicitar cupo (UC1)}
\label{tab:uc1}
\small
\begin{tabularx}{\textwidth}{@{}l X@{}}
\toprule
\rowcolor{azulClaro}
\textbf{Elemento} & \textbf{Descripción} \\
\midrule
Actor principal & Estudiante \\
Actores secundarios & Sistema de notificaciones, Secretaría \\
Precondición & El estudiante está autenticado y el periodo de solicitudes está abierto. \\
Flujo principal &
  1. El estudiante consulta el catálogo de materias y su disponibilidad. \newline
  2. Selecciona la materia y confirma la solicitud. \newline
  3. El sistema valida prerrequisitos, cupo y conflicto de horario. \newline
  4. Si hay cupo, se reserva y se notifica la aprobación. \newline
  5. Si no hay cupo, se registra en lista de espera y se notifica la posición. \\
Flujos alternativos &
  A1: No cumple prerrequisitos $\to$ rechazo con motivo explícito. \newline
  A2: Conflicto de horario $\to$ se sugiere horario alternativo. \newline
  A3: Cupo liberado por no confirmación $\to$ se notifica al siguiente en espera. \\
Postcondición & La solicitud queda registrada con estado y trazabilidad completa. \\
Requerimientos & RF-05 a RF-10, RN-01 a RN-04 \\
\bottomrule
\end{tabularx}
\end{table}

\begin{table}[H]
\centering
\caption{Especificación del caso de uso: Gestión de queja o sugerencia (UC8)}
\label{tab:uc8}
\small
\begin{tabularx}{\textwidth}{@{}l X@{}}
\toprule
\rowcolor{naranjaClaro}
\textbf{Elemento} & \textbf{Descripción} \\
\midrule
Actor principal & Estudiante o Docente \\
Actores secundarios & Secretaría, Docente involucrado, Departamento \\
Precondición & El usuario está autenticado. \\
Flujo principal &
  1. El usuario registra la queja o sugerencia con categoría y descripción. \newline
  2. El sistema clasifica automáticamente urgencia y área responsable. \newline
  3. La secretaría revisa, asigna responsable y plazo. \newline
  4. La secretaría se comunica formalmente con el docente si aplica. \newline
  5. Se registra la acción, se cierra el caso y se solicita satisfacción. \\
Flujos alternativos &
  A1: Caso crítico sin atención en 3 días $\to$ escalamiento a dirección. \newline
  A2: El denunciante solicita anonimato $\to$ se oculta su identidad en todo el flujo. \\
Postcondición & Caso cerrado con acción registrada, evidencia y encuesta de satisfacción. \\
Requerimientos & RF-35 a RF-42, RN-09, RN-10 \\
\bottomrule
\end{tabularx}
\end{table}

\begin{table}[H]
\centering
\caption{Especificación del caso de uso: Evaluación docente a tiempo (UC10)}
\label{tab:uc10}
\small
\begin{tabularx}{\textwidth}{@{}l X@{}}
\toprule
\rowcolor{verdeClaro}
\textbf{Elemento} & \textbf{Descripción} \\
\midrule
Actor principal & Departamento (Administrador) \\
Actores secundarios & Estudiante, Docente, Sistema de notificaciones \\
Precondición & El periodo académico está configurado con fechas de evaluación. \\
Flujo principal &
  1. El departamento configura periodo, preguntas y fecha límite. \newline
  2. El sistema habilita el formulario en la fecha de inicio. \newline
  3. Se envían recordatorios automáticos a los pendientes. \newline
  4. Al cumplirse la fecha límite se consolidan los resultados. \newline
  5. Se publican resultados anonimizados y se generan alertas de desempeño atípico. \\
Flujos alternativos &
  A1: Participación insuficiente $\to$ se extiende el plazo y se intensifican recordatorios. \newline
  A2: El docente solicita revisión $\to$ se verifica trazabilidad sin revelar identidades. \\
Postcondición & 100\% de evaluaciones aplicadas y resultados publicados en el periodo. \\
Requerimientos & RF-27 a RF-34, RN-05, RN-06 \\
\bottomrule
\end{tabularx}
\end{table}

\section{Diagrama de Clases}

\subsection{Descripción general}
El diagrama de clases define la estructura estática del \ac{SGDIC}. Se utiliza una
clase abstracta \texttt{Usuario} de la que heredan los cuatro roles del sistema.
Las entidades de dominio modelan los procesos problemáticos identificados.

\begin{figure}[H]
\centering
\resizebox{\textwidth}{!}{%
\begin{tikzpicture}[
  entidad/.style={
    draw=azulOscuro, fill=azulClaro, thick, rounded corners=3pt,
    text=azulOscuro, align=center, font=\scriptsize\bfseries,
    minimum width=3.35cm, minimum height=0.9cm, inner sep=4pt
  },
  abstracta/.style={
    draw=verdeUC, fill=verdeClaro, thick, dashed, rounded corners=3pt,
    text=verdeUC!60!black, align=center, font=\scriptsize\itshape\bfseries,
    minimum width=3.9cm, minimum height=0.9cm, inner sep=4pt
  },
  etiqueta/.style={
    font=\tiny, fill=white, inner sep=1.2pt, text=azulOscuro
  },
  modulo/.style={
    draw=grisUC!55, fill=grisClaro!60, rounded corners=5pt, inner sep=8pt
  },
  node distance=1.0cm and 1.0cm
]

% --- Jerarquía de usuarios ---
\node[abstracta] (usuario) at (0,0) {Usuario (abstracta)};
\node[entidad] (est) at (-7.2,-2.0) {Estudiante};
\node[entidad] (doc) at (-2.4,-2.0) {Docente};
\node[entidad] (sec) at (2.4,-2.0) {Secretaría};
\node[entidad] (adm) at (7.2,-2.0) {Administrador};

\draw[generalizacion] (est.north) -- ++(0,0.65) -| (usuario.south west);
\draw[generalizacion] (doc.north) -- ++(0,0.65) -| (usuario.south);
\draw[generalizacion] (sec.north) -- ++(0,0.65) -| (usuario.south);
\draw[generalizacion] (adm.north) -- ++(0,0.65) -| (usuario.south east);

% --- Módulos de dominio ---
\node[entidad] (mat) at (-8.6,-5.0) {Materia};
\node[entidad] (sol) at (-5.4,-5.0) {SolicitudCupo};

\node[entidad] (msg) at (-3.4,-7.5) {Mensaje};
\node[entidad] (que) at (0.0,-7.5) {QuejaSugerencia};

\node[entidad] (fal) at (-3.4,-10.0) {Falta};
\node[entidad] (cal) at (0.0,-10.0) {Calificacion};

\node[entidad] (cue) at (5.2,-5.0) {CuestionarioHorario};
\node[entidad] (eva) at (8.8,-5.0) {EvaluacionDocente};

\node[entidad] (rep) at (3.8,-9.0) {Reporte};
\node[entidad] (ana) at (7.2,-9.0) {AnaliticaDatos};
\node[entidad] (rec) at (7.2,-11.55) {Recomendacion};

\begin{scope}[on background layer]
\node[modulo, fit={(mat) (sol)}] {};
\node[modulo, fit={(msg) (que) (fal) (cal)}] {};
\node[modulo, fit={(cue) (eva) (rep) (ana) (rec)}] {};
\end{scope}

\node[font=\tiny\bfseries, text=grisUC] at (-7.0,-4.15) {Gestión académica};
\node[font=\tiny\bfseries, text=grisUC] at (-1.7,-6.65) {Comunicación y seguimiento};
\node[font=\tiny\bfseries, text=grisUC] at (6.5,-4.15) {Evaluación y analítica};

% --- Relaciones principales con rutas ortogonales ---
\draw[thick, azulOscuro] (est.south) |- node[etiqueta, near end, above]{realiza} (sol.west);
\draw[thick, azulOscuro] (mat.east) -- node[etiqueta, above]{recibe} (sol.west);
\draw[thick, azulOscuro] (doc.south) |- node[etiqueta, near end, above]{imparte} (mat.east);

\draw[thick, azulOscuro] (est.south) |- node[etiqueta, near end, above]{envía} (msg.west);
\draw[thick, azulOscuro] (doc.south) -- node[etiqueta, right]{envía} (msg.north);
\draw[thick, azulOscuro] (est.south) |- node[etiqueta, near end, above]{reporta} (que.west);
\draw[thick, azulOscuro] (sec.south) |- node[etiqueta, near end, above]{gestiona} (que.east);

\draw[thick, azulOscuro] (doc.south) |- node[etiqueta, near end, above]{registra} (fal.west);
\draw[thick, azulOscuro] (est.south) |- node[etiqueta, pos=0.74, above]{acumula} (fal.west);
\draw[thick, azulOscuro] (est.south) |- node[etiqueta, near end, above]{califica} (cal.west);
\draw[thick, azulOscuro] (doc.south) |- node[etiqueta, near end, above]{recibe} (cal.west);

\draw[thick, azulOscuro] (est.south) |- node[etiqueta, near end, above]{responde} (cue.west);
\draw[thick, azulOscuro] (adm.south) |- node[etiqueta, near end, above]{configura} (eva.east);
\draw[thick, azulOscuro] (sec.south) |- node[etiqueta, near end, above]{genera} (rep.west);
\draw[thick, azulOscuro] (adm.south) |- node[etiqueta, near end, above]{ejecuta} (ana.east);
\draw[thick, azulOscuro] (ana.south) -- node[etiqueta, right]{sugiere} (rec.north);

\end{tikzpicture}%
}
\caption{Diagrama de clases (visión de relaciones principales) del \ac{SGDIC}}
\label{fig:clases-relaciones}
\end{figure}

\subsection{Diagrama de clases detallado}
La figura~\ref{fig:clases-detalle} presenta las clases principales con sus
atributos y métodos en notación \ac{UML} completa (visibilidad: $-$ privado,
$+$ público).

\begin{figure}[H]
\centering
\resizebox{\textwidth}{!}{%
\begin{tikzpicture}[
  clasedet/.style={
    draw=azulOscuro, fill=white, thick,
    text=black, align=left, font=\ttfamily\tiny, inner sep=3pt
  }
]

% ============ Usuario (abstracta) ============
\node[clasedet] (usuario) at (0,0) {
\begin{tabular}{@{}l@{}}
\multicolumn{1}{@{}c@{}}{\textbf{\textsf{Usuario}} \textit{\{abstract\}}}\\
\hline
$-$ id: String\\
$-$ nombre: String\\
$-$ email: String\\
$-$ passwordHash: String\\
$-$ activo: Boolean\\
\hline
$+$ autenticar(): Boolean\\
$+$ actualizarPerfil(): void\\
\end{tabular}
};

% ============ Estudiante ============
\node[clasedet] (est) at (-9.0,-4.4) {
\begin{tabular}{@{}l@{}}
\multicolumn{1}{@{}c@{}}{\textbf{\textsf{Estudiante}}}\\
\hline
$-$ codigo: String\\
$-$ semestre: Integer\\
$-$ promedio: Double\\
\hline
$+$ solicitarCupos(m): Boolean\\
$+$ enviarMensaje(d,c): Mensaje\\
$+$ calificarDocente(d,p): void\\
$+$ responderCuestionario(): void\\
$+$ presentarQueja(q): void\\
\end{tabular}
};

% ============ Docente ============
\node[clasedet] (doc) at (-3.0,-4.4) {
\begin{tabular}{@{}l@{}}
\multicolumn{1}{@{}c@{}}{\textbf{\textsf{Docente}}}\\
\hline
$-$ codigo: String\\
$-$ especialidad: String\\
$-$ materiasAsignadas: List\\
\hline
$+$ registrarFalta(e,f): Falta\\
$+$ crearEditarMateria(m): void\\
$+$ enviarMensaje(e,c): Mensaje\\
$+$ justificarFalta(f): void\\
\end{tabular}
};

% ============ Secretaría ============
\node[clasedet] (sec) at (3.0,-4.4) {
\begin{tabular}{@{}l@{}}
\multicolumn{1}{@{}c@{}}{\textbf{\textsf{Secretaría}}}\\
\hline
$-$ area: String\\
$-$ casosAsignados: Integer\\
\hline
$+$ gestionarQueja(q): void\\
$+$ comunicarDocente(d,c): void\\
$+$ verificarResultados(): Reporte\\
$+$ generarAnalitica(): Dataset\\
\end{tabular}
};

% ============ Administrador ============
\node[clasedet] (adm) at (9.0,-4.4) {
\begin{tabular}{@{}l@{}}
\multicolumn{1}{@{}c@{}}{\textbf{\textsf{Administrador}}}\\
\hline
$-$ nivelAcceso: Enum\\
\hline
$+$ asignarCupos(): void\\
$+$ configurarEvaluacion(): void\\
$+$ generarSolucion(a): Recomendacion\\
$+$ publicarResultados(): void\\
\end{tabular}
};

% ============ Materia ============
\node[clasedet] (mat) at (-9.0,-8.8) {
\begin{tabular}{@{}l@{}}
\multicolumn{1}{@{}c@{}}{\textbf{\textsf{Materia}}}\\
\hline
$-$ idMateria: String\\
$-$ codigo: String\\
$-$ creditos: Integer\\
$-$ cupoMaximo: Integer\\
$-$ cuposDisponibles: Integer\\
$-$ horarios: List$<$String$>$\\
\hline
$+$ actualizarCupos(): void\\
$+$ verificarDisponibilidad(): Boolean\\
\end{tabular}
};

% ============ QuejaSugerencia ============
\node[clasedet] (que) at (-3.0,-8.8) {
\begin{tabular}{@{}l@{}}
\multicolumn{1}{@{}c@{}}{\textbf{\textsf{QuejaSugerencia}}}\\
\hline
$-$ id: String\\
$-$ tipo: \{QUEJA, SUGERENCIA\}\\
$-$ estado: \{ABIERTO, EN\_PROCESO, RESUELTO\}\\
$-$ prioridad: \{BAJA, MEDIA, ALTA\}\\
$-$ responsable: String\\
$-$ fechaCierre: Date\\
\hline
$+$ asignarSecretaria(): void\\
$+$ generarRespuesta(): void\\
$+$ cerrar(evidencia): void\\
\end{tabular}
};

% ============ EvaluacionDocente ============
\node[clasedet] (eva) at (3.0,-8.8) {
\begin{tabular}{@{}l@{}}
\multicolumn{1}{@{}c@{}}{\textbf{\textsf{EvaluacionDocente}}}\\
\hline
$-$ id: String\\
$-$ periodo: String\\
$-$ fechaInicio: Date\\
$-$ fechaLimite: Date\\
$-$ completadas: Integer\\
$-$ total: Integer\\
\hline
$+$ habilitar(): void\\
$+$ enviarRecordatorio(): void\\
$+$ consolidar(): Reporte\\
$+$ verificarCumplimiento(): Boolean\\
\end{tabular}
};

% ============ AnaliticaDatos ============
\node[clasedet] (ana) at (9.0,-8.8) {
\begin{tabular}{@{}l@{}}
\multicolumn{1}{@{}c@{}}{\textbf{\textsf{AnaliticaDatos}}}\\
\hline
$-$ dataset: Dataset\\
$-$ metricas: Map\\
\hline
$+$ calcularKPI(): Map\\
$+$ detectarAnomalias(): List\\
$+$ predecirDemanda(): Modelo\\
$+$ priorizarRecomendaciones(): List\\
\end{tabular}
};

% --- Relaciones de generalización ---
\draw[generalizacion] (est) -- (usuario);
\draw[generalizacion] (doc) -- (usuario);
\draw[generalizacion] (sec) -- (usuario);
\draw[generalizacion] (adm) -- (usuario);

% --- Asociaciones ---
\draw[thick, azulOscuro] (adm) -- (mat) node[midway, above, font=\tiny, text=azulOscuro]{1 .. *};
\draw[thick, azulOscuro] (est) -- (mat);
\draw[thick, azulOscuro] (sec) -- (que) node[midway, above, font=\tiny, text=azulOscuro]{0 .. *};
\draw[thick, azulOscuro] (adm) -- (eva);
\draw[thick, azulOscuro] (adm) -- (ana);
\draw[thick, azulOscuro] (doc) -- (eva) ;

\end{tikzpicture}%
}
\caption{Diagrama de clases detallado con atributos y métodos}
\label{fig:clases-detalle}
\end{figure}

\section{Diagrama de Secuencia: Evaluación Docente}

El diagrama de secuencia modela la interacción temporal del proceso de
evaluación docente, crítico para garantizar que se complete dentro del periodo
académico.

\begin{figure}[H]
\centering
\resizebox{\textwidth}{!}{%
\begin{tikzpicture}[
  obj/.style={
    draw=azulOscuro, fill=azulClaro, thick, rounded corners=2pt,
    text=azulOscuro, align=center, font=\scriptsize\bfseries,
    minimum width=2.6cm, minimum height=0.85cm, inner sep=3pt
  },
  linea/.style={draw=grisUC, dashed, thick},
  men/.style={->, >=Stealth, thick, color=azulOscuro},
  ret/.style={->, >=Stealth, thick, dashed, color=verdeUC},
  self/.style={->, >=Stealth, thick, color=moradoUC}
]

% --- Objetos (líneas de vida) ---
\node[obj] (adm)   at (0,0)    {Administrador};
\node[obj] (sys)   at (3.6,0)  {SGDIC};
\node[obj] (est)   at (7.2,0)  {Estudiante};
\node[obj] (doc)   at (10.8,0) {Docente};
\node[obj] (sec)   at (14.4,0) {Secretaría};

% --- Líneas de vida ---
\draw[linea] (adm) -- (0,-13.0);
\draw[linea] (sys) -- (3.6,-13.0);
\draw[linea] (est) -- (7.2,-13.0);
\draw[linea] (doc) -- (10.8,-13.0);
\draw[linea] (sec) -- (14.4,-13.0);

% --- Mensajes ---
\draw[men] (0,-0.9) -- node[above, font=\tiny]{1. configurarEvaluacion(periodo, fechaLimite)} (3.6,-0.9);
\draw[self] (3.6,-1.4) -- node[right, font=\tiny]{2. crear formulario} (3.6,-1.4);
\draw[men] (3.6,-2.1) -- node[above, font=\tiny]{3. habilitarEvaluacion()} (7.2,-2.1);
\draw[men] (7.2,-2.8) -- node[above, font=\tiny]{4. mostrar formulario anónimo} (3.6,-2.8);
\draw[men] (7.2,-3.6) -- node[above, font=\tiny]{5. responderEvaluacion(puntajes)} (3.6,-3.6);
\draw[self] (3.6,-4.2) -- node[right, font=\tiny]{6. validar ventana de tiempo} (3.6,-4.2);
\draw[ret] (3.6,-4.9) -- node[above, font=\tiny]{7. confirmarRegistro()} (7.2,-4.9);

% --- Recordatorios ---
\draw[self] (3.6,-5.7) -- node[right, font=\tiny]{8. enviarRecordatorio()} (3.6,-5.7);
\draw[men] (3.6,-6.4) -- node[above, font=\tiny]{9. notificar pendientes} (7.2,-6.4);
\draw[men] (3.6,-7.2) -- node[above, font=\tiny]{10. alertarNoCompletadas()} (14.4,-7.2);

% --- Cierre y consolidación ---
\draw[men] (14.4,-8.0) -- node[above, font=\tiny]{11. verificarCumplimiento()} (3.6,-8.0);
\draw[self] (3.6,-8.7) -- node[right, font=\tiny]{12. consolidar resultados} (3.6,-8.7);
\draw[men] (3.6,-9.4) -- node[above, font=\tiny]{13. publicarResultados(docente)} (10.8,-9.4);
\draw[men] (3.6,-10.1) -- node[above, font=\tiny]{14. generarReporte(secretaría)} (14.4,-10.1);
\draw[men] (3.6,-10.8) -- node[above, font=\tiny]{15. analizarDesempeñoAtípico()} (0,-10.8);
\draw[ret] (0,-11.5) -- node[above, font=\tiny]{16. plan de acompañamiento} (3.6,-11.5);

% --- Nota de anonimato ---
\node[draw=moradoUC, fill=moradoClaro, rounded corners=2pt, font=\tiny,
      align=left, inner sep=3pt, text width=4.0cm] at (7.2,-12.4) {
  \textbf{Restricción:} el docente nunca recibe la identidad del evaluador;
  solo resultados agregados con mínimo 5 respuestas.
};

\end{tikzpicture}%
}
\caption{Diagrama de secuencia del proceso de evaluación docente}
\label{fig:secuencia-evaluacion}
\end{figure}

\section{Diagrama de Actividades: Gestión de Quejas y Sugerencias}

El diagrama de actividades describe el flujo completo de una queja o sugerencia,
desde su registro hasta su cierre verificado, incluyendo el escalamiento
automático de casos críticos.

\begin{figure}[H]
\centering
\resizebox{0.95\textwidth}{!}{%
\begin{tikzpicture}[node distance=0.5cm, every node/.style={font=\scriptsize}]

% --- Swimlane: Estudiante ---
\node[terminal] (ini) {Inicio};
\node[flujo, below=of ini] (a1) {Registrar queja o\\sugerencia};
\node[flujo, below=2.5cm of a1] (a2) {Recibir\\notificación de estado};
\node[flujo, below=of a2] (a3) {Responder encuesta\\de satisfacción};
\node[terminal, below=of a3] (fin) {Fin};

% --- Flujo del sistema ---
\node[flujo, right=2.7cm of a1] (s1) {Clasificar categoría,\\urgencia y área};
\node[decision, right=2.7cm of s1, xshift=0.5cm] (d1) {¿Caso\\crítico?};
\node[flujo, right=2.7cm of d1, xshift=0.5cm] (s2) {Escalar a\\dirección};

% --- Flujo de secretaría ---
\node[flujo, below=1.0cm of s1] (sec1) {Asignar responsable\\y plazo};
\node[decision, below=1.6cm of sec1] (d2) {¿Vencido\\el plazo?};
\node[flujo, right=2.7cm of d2] (sec2) {Reasignar y\\alertar};
\node[flujo, below=1.0cm of d2] (sec3) {Contactar al\\docente involucrado};
\node[flujo, below=of sec3] (sec4) {Registrar acción\\y evidencia};
\node[decision, below=1.6cm of sec4] (d3) {¿Acción\\suficiente?};
\node[flujo, right=2.7cm of d3] (sec5) {Escalar a\\dirección};
\node[flujo, below=1.0cm of d3] (sec6) {Cerrar caso y\\notificar};

% --- Conexiones ---
\draw[flecha] (ini) -- (a1);
\draw[flecha] (a1) -- (s1);
\draw[flecha] (s1) -- (d1);
\draw[flecha] (d1) -- node[above, font=\tiny]{Sí} (s2);
\draw[flecha] (d1) -- node[right, font=\tiny]{No} (sec1);
\draw[flecha] (s2) |- (sec3);
\draw[flecha] (sec1) -- (d2);
\draw[flecha] (d2) -- node[above, font=\tiny]{Sí} (sec2);
\draw[flecha] (sec2) |- (sec3);
\draw[flecha] (d2) -- node[right, font=\tiny]{No} (sec3);
\draw[flecha] (sec3) -- (sec4);
\draw[flecha] (sec4) -- (d3);
\draw[flecha] (d3) -- node[above, font=\tiny]{No} (sec5);
\draw[flecha] (sec5) |- (sec4);
\draw[flecha] (d3) -- node[right, font=\tiny]{Sí} (sec6);
\draw[flecha] (sec6) -- (a2);
\draw[flecha] (a2) -- (a3);
\draw[flecha] (a3) -- (fin);

\end{tikzpicture}%
}
\caption{Diagrama de actividades del flujo de quejas y sugerencias}
\label{fig:actividades-quejas}
\end{figure}

\section{Diagrama de Estados: Ciclo de Vida de una Queja}

\begin{figure}[H]
\centering
\begin{tikzpicture}[node distance=1.1cm and 1.3cm, every node/.style={font=\scriptsize}]

\node[terminal] (ini) {Inicio};
\node[flujo, below=of ini] (registrada) {REGISTRADA};
\node[flujo, below=of registrada] (clasificada) {CLASIFICADA};
\node[flujo, below left=of clasificada, xshift=-0.6cm] (asignada) {ASIGNADA};
\node[flujo, below right=of clasificada, xshift=0.6cm] (escalada) {ESCALADA};
\node[flujo, below=1.6cm of clasificada] (enproceso) {EN PROCESO};
\node[flujo, below=of enproceso] (atendida) {ATENDIDA};
\node[flujo, below=of atendida] (cerrada) {CERRADA};
\node[flujo, right=2.6cm of cerrada] (reabierta) {REABIERTA};
\node[terminal, below=of cerrada] (fin) {Fin};

\draw[flecha] (ini) -- (registrada);
\draw[flecha] (registrada) -- (clasificada);
\draw[flecha] (clasificada) -- (asignada);
\draw[flecha] (clasificada) -- (escalada);
\draw[flecha] (asignada) -- (enproceso);
\draw[flecha] (escalada) |- (enproceso);
\draw[flecha] (enproceso) -- (atendida);
\draw[flecha] (atendida) -- (cerrada);
\draw[flecha] (cerrada) -- (fin);
\draw[flecha] (cerrada) -- node[above, font=\tiny]{insatisfacción} (reabierta);
\draw[flecha] (reabierta) |- (enproceso);

% --- Transiciones etiquetadas ---
\node[right=0.15cm of registrada, font=\tiny, text=grisUC, align=left] {por el usuario};
\node[right=0.15cm of clasificada, font=\tiny, text=grisUC, align=left] {automática};
\node[left=0.15cm of asignada, font=\tiny, text=grisUC, align=left] {secretaría};
\node[right=0.15cm of escalada, font=\tiny, text=grisUC, align=left] {crítico $>$ 3 días};
\node[left=0.15cm of atendida, font=\tiny, text=grisUC, align=left] {acción registrada};

\end{tikzpicture}
\caption{Diagrama de estados del ciclo de vida de una queja o sugerencia}
\label{fig:estados-queja}
\end{figure}

\section{Diagrama de Componentes y Arquitectura}

\begin{figure}[H]
\centering
\begin{tikzpicture}[
  comp/.style={
    draw=azulOscuro, fill=azulClaro, thick, rounded corners=3pt,
    text=azulOscuro, align=center, font=\scriptsize\bfseries,
    minimum width=3.0cm, minimum height=1.0cm, inner sep=4pt
  },
  infra/.style={
    draw=verdeUC, fill=verdeClaro, thick, rounded corners=3pt,
    text=verdeUC!60!black, align=center, font=\scriptsize\bfseries,
    minimum width=3.0cm, minimum height=1.0cm, inner sep=4pt
  },
  node distance=0.6cm and 1.0cm
]

% --- Capa de presentación ---
\node[comp] (web) {Portal Web\\(Estudiante/Docente)};
\node[comp, right=of web] (admin) {Consola\\Administrativa};
\node[comp, right=of admin] (movil) {Vista Móvil\\(Responsiva)};

% --- Capa de servicios ---
\node[comp, below=1.0cm of web] (srv1) {Servicio de\\Cupos y Materias};
\node[comp, below=1.0cm of admin] (srv2) {Servicio de\\Comunicación};
\node[comp, below=1.0cm of movil] (srv3) {Servicio de\\Evaluación y Quejas};

% --- Capa analítica ---
\node[infra, below=1.0cm of srv2] (etl) {Motor \ac{ETL}};
\node[infra, right=of etl] (dwh) {Repositorio\\Analítico};
\node[infra, left=of etl] (ml) {Motor de\\Soluciones Rápidas};

% --- Integraciones ---
\node[infra, below=1.0cm of etl] (api) {\ac{API} \ac{REST} Institucional};
\node[infra, left=of api] (notif) {Servicio de\\Notificaciones};
\node[infra, right=of api] (auth) {Autenticación\\Institucional};

% --- Conexiones ---
\draw[flecha] (web) -- (srv1);
\draw[flecha] (admin) -- (srv2);
\draw[flecha] (movil) -- (srv3);
\draw[flecha] (srv1) -- (etl);
\draw[flecha] (srv2) -- (etl);
\draw[flecha] (srv3) -- (etl);
\draw[flecha] (etl) -- (dwh);
\draw[flecha] (dwh) -- (ml);
\draw[flecha] (ml) -- (srv1);
\draw[flecha] (api) -- (etl);
\draw[flecha] (notif) -- (srv2);
\draw[flecha] (auth) -- (srv3);

\end{tikzpicture}
\caption{Arquitectura de componentes del \ac{SGDIC}}
\label{fig:componentes}
\end{figure}

\begin{table}[H]
\centering
\caption{Descripción de los componentes de la arquitectura}
\label{tab:componentes}
\small
\begin{tabularx}{\textwidth}{@{}l X@{}}
\toprule
\rowcolor{azulClaro}
\textbf{Componente} & \textbf{Responsabilidad} \\
\midrule
Portal Web & Interfaz principal para estudiantes y docentes (gestión académica). \\
Consola Administrativa & Configuración, verificación y tableros para secretaría y departamento. \\
Vista Móvil & Acceso responsivo optimizado para notificaciones y consultas rápidas. \\
Servicio de Cupos y Materias & Lógica de disponibilidad, priorización y ciclo de vida de materias. \\
Servicio de Comunicación & Mensajería trazable y notificaciones automáticas. \\
Servicio de Evaluación y Quejas & Ventanas de evaluación, cuestionarios y gestión de casos. \\
Motor \ac{ETL} & Extracción, limpieza, transformación y carga de datos al repositorio. \\
Repositorio Analítico & Almacén consolidado para cálculo de \ac{KPI} y modelos. \\
Motor de Soluciones Rápidas & Detección de anomalías, predicciones y generación de recomendaciones. \\
\ac{API} \ac{REST} Institucional & Integración con matrícula, catálogo oficial y autenticación. \\
Servicio de Notificaciones & Despacho de correos, mensajes internos y recordatorios. \\
Autenticación Institucional & Validación de credenciales y roles (\ac{RBAC}). \\
\bottomrule
\end{tabularx}
\end{table}

\section{Resumen de Diagramas}

\begin{table}[H]
\centering
\caption{Inventario de diagramas \ac{UML} del documento}
\label{tab:inventario-uml}
\small
\begin{tabularx}{\textwidth}{@{}p{3.1cm} p{2.2cm} X@{}}
\toprule
\rowcolor{verdeClaro}
\textbf{Diagrama} & \textbf{Figura} & \textbf{Propósito} \\
\midrule
Casos de uso & \ref{fig:casos-uso} & Funcionalidades y actores del sistema. \\
Clases (relaciones) & \ref{fig:clases-relaciones} & Estructura estática y asociaciones. \\
Clases (detalle) & \ref{fig:clases-detalle} & Atributos y métodos por clase. \\
Secuencia & \ref{fig:secuencia-evaluacion} & Interacción temporal de la evaluación docente. \\
Actividades & \ref{fig:actividades-quejas} & Flujo completo de quejas y sugerencias. \\
Estados & \ref{fig:estados-queja} & Ciclo de vida de una queja o sugerencia. \\
Componentes & \ref{fig:componentes} & Arquitectura de la solución. \\
\bottomrule
\end{tabularx}
\end{table}
% =====================================================================
% CAPÍTULO 6: ALGORITMO DE SOLUCIONES RÁPIDAS Y ANALÍTICA DE DATOS
% =====================================================================
\chapter{Algoritmo de Soluciones Rápidas y Analítica de Datos}
\label{ch:algoritmo}

\section{Definición de Solución Rápida}

\begin{cajaAlgoritmo}[Definición]
Una \textbf{solución rápida} es una recomendación accionable, priorizada y
trazable que el motor analítico del \ac{SGDIC} genera para resolver una
situación detectada (saturación de cupos, queja recurrente, acumulación de
inasistencia, desempeño docente atípico o conflicto de horarios), indicando
\textbf{qué hacer, quién es el responsable, en qué plazo y con qué impacto
esperado}.
\end{cajaAlgoritmo}

\section{Arquitectura del Motor Analítico}

\begin{figure}[H]
\centering
\resizebox{0.92\textwidth}{!}{%
\begin{tikzpicture}[node distance=0.6cm and 0.8cm, every node/.style={font=\scriptsize}]

\node[terminal] (datos) {Datos operativos del \ac{SGDIC}};
\node[flujo, below=of datos] (etl) {1. Extracción y limpieza (\ac{ETL})};
\node[flujo, below=of etl] (feat) {2. Ingeniería de variables};
\node[flujo, below=of feat] (modelos) {3. Modelos analíticos};

\node[flujo, below left=1.0cm and 0.5cm of modelos] (m1) {Clasificación\\de quejas};
\node[flujo, below=1.0cm of modelos] (m2) {Predicción de\\demanda de cupos};
\node[flujo, below right=1.0cm and 0.5cm of modelos] (m3) {Preferencia de\\horarios};

\node[flujo, below=3.0cm of modelos] (prior) {4. Motor de priorización\\\scriptsize (impacto $\times$ urgencia $\div$ esfuerzo)};
\node[flujo, below=of prior] (rec) {5. Generación de\\soluciones rápidas};
\node[decision, below=of rec] (val) {¿Recomendación\\validada?};
\node[flujo, right=1.6cm of val] (ajuste) {Ajustar\\parámetros};
\node[terminal, below=of val] (tablero) {6. Tablero de\\verificación};

\draw[flecha] (datos) -- (etl);
\draw[flecha] (etl) -- (feat);
\draw[flecha] (feat) -- (modelos);
\draw[flecha] (modelos) -- (m1);
\draw[flecha] (modelos) -- (m2);
\draw[flecha] (modelos) -- (m3);
\draw[flecha] (m1) -- (prior);
\draw[flecha] (m2) -- (prior);
\draw[flecha] (m3) -- (prior);
\draw[flecha] (prior) -- (rec);
\draw[flecha] (rec) -- (val);
\draw[flecha] (val) -- node[right, font=\tiny]{Sí} (tablero);
\draw[flecha] (val) -- node[above, font=\tiny]{No} (ajuste);
\draw[flecha] (ajuste) |- (modelos);
\draw[flecha] (tablero.east) -- ++(1.6,0) |- node[near start, right, font=\tiny]{retroalimentación} (feat.east);

\end{tikzpicture}%
}
\caption{Arquitectura del motor analítico y de soluciones rápidas}
\label{fig:motor-analitico}
\end{figure}

\section{Descripción del Algoritmo}

\subsection{Vista general}
El algoritmo consta de seis etapas que transforman datos operativos en
recomendaciones priorizadas.

\begin{table}[H]
\centering
\caption{Etapas del algoritmo de soluciones rápidas}
\label{tab:etapas-algoritmo}
\small
\begin{tabularx}{\textwidth}{@{}c l X@{}}
\toprule
\rowcolor{moradoClaro}
\textbf{Etapa} & \textbf{Nombre} & \textbf{Operación} \\
\midrule
1 & Ingesta y limpieza &
    Consolidación de datos de cupos, mensajes, faltas, evaluaciones, quejas y
    cuestionarios; eliminación de duplicados y normalización. \\
2 & Ingeniería de variables &
    Cálculo de variables derivadas: tasa de ocupación, tiempo medio de
    respuesta, índice de inasistencia, sentimiento de quejas, frecuencia
    histórica de preferencia horaria. \\
3 & Modelado &
    Ejecución de modelos de clasificación, regresión y probabilidad según el
    tipo de problema detectado. \\
4 & Detección de anomalías &
    Identificación de desviaciones significativas respecto al comportamiento
    histórico esperado. \\
5 & Priorización &
    Ordenamiento de hallazgos según impacto potencial, urgencia y esfuerzo de
    implementación. \\
6 & Generación y despliegue &
    Construcción del texto de la recomendación, asignación de responsable,
    plazo sugerido y publicación en el tablero. \\
\bottomrule
\end{tabularx}
\end{table}

\subsection{Pseudocódigo del algoritmo}

\begin{lstlisting}[language=Python, caption={Algoritmo de generación de soluciones rápidas}, label={lst:algoritmo}]
ALGORITMO GenerarSolucionesRapidas(periodo):

    # Etapa 1: ingesta y limpieza
    D <- ExtraerDatos(periodo)
    D <- LimpiarYNormalizar(D)
    D <- Anonimizar(D)

    # Etapa 2: ingenieria de variables
    V <- CalcularVariables(D)

    # Etapa 3: modelado
    M_quejas  <- ModeloClasificacionQuejas(V.quejas)
    M_demanda <- ModeloPrediccionCupos(V.cupos)
    M_horario <- ModeloProbabilidadHorario(V.cuestionarios)
    M_desemp  <- ModeloDesempenoDocente(V.evaluaciones)

    # Etapa 4: deteccion de anomalias
    A <- []
    PARA CADA materia m EN V.materias HACER
        SI tasa_ocupacion(m) > 0.90 ENTONCES
            A.agregar(Anomalia("SATURACION_CUPOS", m))
        FIN SI
    FIN PARA

    PARA CADA docente d EN V.docentes HACER
        SI puntaje(d) < umbral_bajo(d) ENTONCES
            A.agregar(Anomalia("DESEMPENO_ATIPICO", d))
        FIN SI
    FIN PARA

    PARA CADA patron p EN DetectarPatrones(V.quejas) HACER
        SI p.frecuencia > umbral_frecuencia ENTONCES
            A.agregar(Anomalia("QUEJA_RECURRENTE", p))
        FIN SI
    FIN PARA

    # Etapa 5: priorizacion
    PARA CADA anomalia a EN A HACER
        a.esfuerzo <- EstimarEsfuerzo(a)
        a.score    <- (EstimarImpacto(a) * EstimarUrgencia(a))
                      / (a.esfuerzo + EPSILON)
    FIN PARA
    A <- OrdenarDescendente(A, por=a.score)

    # Etapa 6: generacion de recomendaciones
    R <- []
    PARA CADA anomalia a EN A HACER
        r <- ConstruirRecomendacion(a)
        r.responsable <- AsignarResponsable(a)
        r.plazo       <- CalcularPlazo(EstimarUrgencia(a))
        r.confianza   <- CalcularConfianza(a)
        r.estado      <- "PENDIENTE"
        R.agregar(r)
    FIN PARA

    PublicarEnTablero(R)
    RETORNAR R

FIN ALGORITMO
\end{lstlisting}

\subsection{Fórmula de priorización}

La puntuación de prioridad de cada recomendación se calcula como:

\begin{equation}
\text{Score} = \frac{I \cdot U}{E + \varepsilon}
\label{eq:score}
\end{equation}

donde $I \in [1,5]$ es el impacto estimado, $U \in [1,5]$ la urgencia,
$E \in [1,5]$ el esfuerzo de implementación y $\varepsilon$ un valor pequeño
que evita la división por cero. Este criterio favorece las acciones de alto
impacto y alta urgencia con bajo esfuerzo, es decir, las soluciones rápidas
propiamente dichas.

\subsection{Modelo de preferencia de horarios por probabilidad}

El módulo de cuestionarios de materias y horarios alimenta el modelo
probabilístico de preferencia horaria. Para cada combinación de perfil de
estudiante $p$ y franja horaria $h_i$, se estima la probabilidad:

\begin{equation}
P(h_i \mid p) =
\frac{P(p \mid h_i) \cdot P(h_i)}{P(p)}
= \frac{N_{p,h_i} + \alpha}{N_{h_i} + \alpha \cdot |P|}
\label{eq:prob-horario}
\end{equation}

donde $N_{p,h_i}$ es el número de estudiantes del perfil $p$ que prefirieron la
franja $h_i$, $N_{h_i}$ el total de preferencias de esa franja, $|P|$ el número
de perfiles y $\alpha$ un parámetro de suavizado de Laplace que evita
probabilidades nulas para franjas sin observaciones.

La preferencia esperada total de una franja se calcula como el promedio
ponderado sobre los perfiles de la población objetivo:

\begin{equation}
P(h_i) = \sum_{p \in P} P(h_i \mid p) \cdot w_p
\label{eq:prob-total}
\end{equation}

donde $w_p$ es el peso relativo del perfil $p$ en la matrícula. La asignación de
horario se resuelve como un problema de optimización que maximiza la aceptación
esperada:

\begin{equation}
\max \sum_{i=1}^{n} P(h_i) \cdot x_i
\quad \text{sujeto a} \quad \sum_{i=1}^{n} x_i = 1, \; x_i \in \{0,1\}
\label{eq:optimizacion}
\end{equation}

donde $x_i = 1$ indica que la franja $h_i$ es seleccionada para la materia.

\section{Modelos Analíticos por Tipo de Problema}

\begin{table}[H]
\centering
\caption{Modelos analíticos aplicados por tipo de problema}
\label{tab:modelos}
\small
\begin{tabularx}{\textwidth}{@{}p{3.2cm} p{6.6cm} X@{}}
\toprule
\rowcolor{azulClaro}
\textbf{Problema} & \textbf{Modelo} & \textbf{Salida} \\
\midrule
Clasificación de quejas &
  Clasificación supervisada (Naive Bayes / \ac{KNN}) sobre el texto de la queja &
  Categoría, área responsable y urgencia sugerida \\
Predicción de demanda de cupos &
  Series temporales (suavizado exponencial / ARIMA) por materia y periodo &
  Demanda esperada y riesgo de saturación \\
Preferencia de horarios &
  Probabilidad condicional con suavizado de Laplace &
  $P(h_i \mid p)$ y ranking de franjas preferidas \\
Desempeño docente atípico &
  Control estadístico de procesos ($\mu \pm 2\sigma$) &
  Docentes con puntaje fuera del rango esperado \\
Reincidencia de inasistencia &
  Regresión logística sobre historial de faltas &
  Probabilidad de superar el umbral permitido \\
Tiempo de respuesta &
  Estadística descriptiva y percentiles &
  Materias y actores con tiempos sobre el objetivo \\
Temas recurrentes en quejas &
  Procesamiento de lenguaje natural (frecuencia de términos) &
  Temas emergentes y su evolución temporal \\
\bottomrule
\end{tabularx}
\end{table}

\section{Reglas de Generación de Soluciones Rápidas}

\begin{table}[H]
\centering
\caption{Catálogo de soluciones rápidas por anomalía detectada}
\label{tab:catalogo-soluciones}
\small
\begin{tabularx}{\textwidth}{@{}p{3.2cm} X p{3.0cm}@{}}
\toprule
\rowcolor{verdeClaro}
\textbf{Anomalía} & \textbf{Solución rápida sugerida} & \textbf{Responsable} \\
\midrule
Saturación de cupos ($>90\%$) &
  Abrir una sección adicional o ampliar el cupo máximo; notificar a los
  estudiantes en lista de espera con horario alternativo. & Departamento \\
Demanda proyectada creciente &
  Planificar apertura de sección antes del inicio del periodo de matrícula. &
  Departamento \\
Queja recurrente en una materia &
  Reunión formal con el docente, revisión del plan de asignatura y plan de
  mejoramiento con seguimiento quincenal. & Secretaría + Departamento \\
Queja crítica sin atender ($>3$ días) &
  Escalamiento inmediato a dirección y notificación al responsable con plazo
  de 24 horas. & Dirección \\
Desempeño docente atípico &
  Activar plan de acompañamiento pedagógico y observación de clase. &
  Departamento \\
Inasistencia acumulada $>20\%$ &
  Alerta al estudiante con plan de recuperación y citación a tutoría. &
  Docente + Secretaría \\
Evaluación docente incompleta &
  Recordatorio intensivo y habilitación de tiempo protegido en clase. &
  Docente + Sistema \\
Conflicto de horarios frecuente &
  Reubicar la materia a la franja con mayor $P(h_i)$ estimada. & Departamento \\
Mensajes sin respuesta $>48$ h &
  Recordatorio automático al destinatario y copia a la secretaría. & Sistema \\
Tiempo de respuesta a cupos elevado &
  Redistribuir la carga de revisión entre el personal de secretaría. &
  Secretaría \\
\bottomrule
\end{tabularx}
\end{table}

\section{Tablero de Verificación de Resultados}

\begin{figure}[H]
\centering
\begin{tikzpicture}[node distance=0.45cm and 0.5cm, every node/.style={font=\scriptsize}]

\node[flujo, minimum width=3.4cm] (tab) at (0,0) {Tablero de\\Verificación};

\node[flujo, below left=1.0cm and 0.3cm of tab] (t1) {Resultados\\académicos};
\node[flujo, below=1.0cm of tab] (t2) {Quejas y\\sugerencias};
\node[flujo, below right=1.0cm and 0.3cm of tab] (t3) {Evaluación\\docente};

\node[flujo, below=1.0cm of t1] (v1) {Ocupación de cupos\\Tiempo de respuesta};
\node[flujo, below=1.0cm of t2] (v2) {Tasa de cierre\\Tiempo por caso};
\node[flujo, below=1.0cm of t3] (v3) {Participación\\Puntaje promedio};

\node[decision, below=1.0cm of v2] (dec) {¿Meta\\alcanzada?};
\node[terminal, left=1.6cm of dec] (ok) {Reportar\\logro};
\node[flujo, right=1.6cm of dec] (no) {Generar acción\\correctiva};

\draw[flecha] (tab) -- (t1);
\draw[flecha] (tab) -- (t2);
\draw[flecha] (tab) -- (t3);
\draw[flecha] (t1) -- (v1);
\draw[flecha] (t2) -- (v2);
\draw[flecha] (t3) -- (v3);
\draw[flecha] (v1) |- (dec);
\draw[flecha] (v3) |- (dec);
\draw[flecha] (v2) -- (dec);
\draw[flecha] (dec) -- node[above, font=\tiny]{Sí} (ok);
\draw[flecha] (dec) -- node[above, font=\tiny]{No} (no);
\draw[flecha] (no) |- (tab);

\end{tikzpicture}
\caption{Estructura del tablero de verificación de resultados}
\label{fig:tablero}
\end{figure}

\subsection{Elementos del tablero}
\begin{table}[H]
\centering
\caption{Componentes del tablero de verificación}
\label{tab:tablero}
\small
\begin{tabularx}{\textwidth}{@{}p{3.3cm} X p{3.1cm}@{}}
\toprule
\rowcolor{moradoClaro}
\textbf{Sección} & \textbf{Contenido} & \textbf{Usuario} \\
\midrule
Panel de cupos & Ocupación por materia, lista de espera, tiempo de respuesta. & Secretaría, Departamento \\
Panel de quejas & Casos abiertos, vencidos, por categoría y responsable. & Secretaría, Dirección \\
Panel de evaluación & Participación por materia, puntajes agregados, alertas. & Departamento, Docente \\
Panel de comunicación & Mensajes sin respuesta, tiempos promedio, escalamientos. & Secretaría \\
Panel de faltas & Estudiantes en riesgo, alertas por umbral, justificaciones pendientes. & Docente, Secretaría \\
Panel de soluciones & Recomendaciones activas, aplicadas y su efecto medido. & Departamento \\
\bottomrule
\end{tabularx}
\end{table}

\section{Métricas de Evaluación del Algoritmo}

\begin{table}[H]
\centering
\caption{Métricas de desempeño del motor analítico}
\label{tab:metricas-algoritmo}
\small
\begin{tabularx}{\textwidth}{@{}p{3.6cm} X c@{}}
\toprule
\rowcolor{azulClaro}
\textbf{Métrica} & \textbf{Definición} & \textbf{Meta} \\
\midrule
Exactitud & $\frac{TP+TN}{TP+TN+FP+FN}$ & $> 0.85$ \\
Precisión & $\frac{TP}{TP+FP}$ & $> 0.80$ \\
Exhaustividad & $\frac{TP}{TP+FN}$ & $> 0.80$ \\
$F_1$ & Media armónica de precisión y exhaustividad & $> 0.80$ \\
Error absoluto medio & $\frac{1}{n}\sum |y_i - \hat{y}_i|$ (demanda de cupos) & $< 10\%$ \\
Tasa de aceptación & Recomendaciones aplicadas / generadas & $> 70\%$ \\
Efectividad & Mejora medida tras la aplicación & $> 60\%$ \\
Cobertura & Anomalías detectadas / anomalías reales & $> 85\%$ \\
\bottomrule
\end{tabularx}
\end{table}

\section{Ejemplo de Aplicación}

\begin{cajaAlgoritmo}[Ejemplo: quejas recurrentes en una materia de programación]
\textbf{Datos observados:} 18 quejas en el periodo correspondiente a la materia
\textit{Programación II}, 14 de ellas con el término ``avance'' y ``ritmo'';
la tasa de inasistencia del grupo es del 24\%; el tiempo medio de respuesta del
docente es de 96 horas.

\vspace{4pt}
\textbf{Hallazgos del motor:}
\begin{itemize}[leftmargin=1.0cm, itemsep=2pt]
  \item Anomalía \texttt{QUEJA\_RECURRENTE} (frecuencia $>$ umbral).
  \item Anomalía \texttt{INASISTENCIA\_ALTA} ($24\% > 20\%$).
  \item Anomalía \texttt{RESPUESTA\_LENTA} ($96$ h $> 48$ h).
\end{itemize}

\vspace{4pt}
\textbf{Priorización:} impacto $I=5$, urgencia $U=4$, esfuerzo $E=2$
$\Rightarrow$ Score $= (5 \times 4)/(2 + 0.01) \approx 9.95$ (prioridad alta).

\vspace{4pt}
\textbf{Solución rápida generada:} reunión formal entre secretaría, docente y
coordinación académica en un plazo de 3 días; revisión del cronograma del plan
de asignatura; tutoría de refuerzo semanal; recordatorios automáticos al docente
sobre mensajes pendientes.
\end{cajaAlgoritmo}
% =====================================================================
% CAPÍTULO 7: PLANIFICACIÓN, CRONOGRAMA Y CONCLUSIONES
% =====================================================================
\chapter{Planificación, Cronograma y Conclusiones}
\label{ch:cronograma}

\section{Cronograma General del Proyecto}

El proyecto se planifica en 24 semanas, distribuidas en ocho fases alineadas con
el Capítulo~\ref{ch:metodologia}.

\begin{figure}[H]
\centering
\resizebox{0.92\textwidth}{!}{%
\begin{ganttchart}[
  hgrid,
  vgrid,
  x unit=0.72cm,
  y unit chart=0.62cm,
  bar height=0.55,
  bar/.append style={fill=azulUC, rounded corners=2pt},
  group/.append style={fill=azulOscuro, draw=azulOscuro},
  milestone/.append style={fill=moradoUC},
  title/.style={fill=azulClaro, draw=azulUC},
  title label font=\tiny\bfseries,
  bar label font=\tiny,
  group label font=\tiny\bfseries,
  milestone label font=\tiny,
  canvas/.append style={draw=azulUC}
]{1}{24}

\gantttitle{Semanas de ejecución}{24} \\
\gantttitlelist{1,2,3,4,5,6,7,8,9,10,11,12,13,14,15,16,17,18,19,20,21,22,23,24}{1} \\

\ganttgroup{Fase 1: Definición}{1}{2} \\
\ganttbar{Diagnóstico del problema}{1}{2} \\
\ganttbar{Formulación de preguntas e hipótesis}{2}{2} \\

\ganttgroup{Fase 2: Revisión de literatura}{2}{5} \\
\ganttbar{Antecedentes y marco teórico}{2}{4} \\
\ganttbar{Normativa y reglas de negocio}{4}{5} \\

\ganttgroup{Fase 3: Diseño metodológico}{6}{8} \\
\ganttbar{Instrumentos y validación}{6}{7} \\
\ganttbar{Definición de muestra}{7}{8} \\

\ganttgroup{Fase 4: Recolección}{8}{12} \\
\ganttbar{Encuestas y entrevistas}{8}{11} \\
\ganttbar{Observación de procesos}{9}{11} \\
\ganttbar{Extracción de datos históricos}{10}{12} \\

\ganttgroup{Fase 5: Análisis}{12}{15} \\
\ganttbar{Codificación cualitativa}{12}{13} \\
\ganttbar{Estadística y preparación del dataset}{13}{15} \\

\ganttgroup{Fase 6: Diseño de la solución}{14}{20} \\
\ganttbar{Requerimientos y modelado UML}{14}{17} \\
\ganttbar{Algoritmo de soluciones rápidas}{17}{20} \\
\ganttbar{Motor analítico y tableros}{18}{20} \\

\ganttgroup{Fase 7: Validación}{19}{22} \\
\ganttbar{Prueba piloto}{19}{21} \\
\ganttbar{Validación con expertos}{21}{22} \\

\ganttgroup{Fase 8: Documentación}{20}{24} \\
\ganttbar{Redacción del informe}{20}{23} \\
\ganttbar{Socialización y cierre}{23}{24} \\

\ganttmilestone{Entrega final}{24}

\end{ganttchart}
}
\caption{Cronograma general del proyecto (24 semanas)}
\label{fig:gantt}
\end{figure}

\section{Plan de Mitigación de Riesgos}

\begin{table}[H]
\centering
\caption{Riesgos, impacto y plan de mitigación}
\label{tab:riesgos}
\small
\begin{tabularx}{\textwidth}{@{}p{1.2cm} X c c X@{}}
\toprule
\rowcolor{naranjaClaro}
\textbf{ID} & \textbf{Riesgo} & \textbf{Prob.} & \textbf{Impacto} & \textbf{Mitigación} \\
\midrule
R-01 & Baja participación en encuestas y evaluación docente. &
  Media & Alto &
  Recordatorios automáticos, tiempo protegido en clase, incentivos
  académicos y reducción de la longitud del instrumento. \\
R-02 & Resistencia al cambio de docentes y personal administrativo. &
  Media & Alto &
  Socialización temprana, capacitación, pilotos con voluntarios y
  visibilización de beneficios. \\
R-03 & Datos históricos incompletos o inconsistentes. &
  Alta & Medio &
  Reglas de limpieza documentadas, imputación controlada y exclusión de
  registros no confiables. \\
R-04 & Percepción de vigilancia sobre la actividad docente. &
  Media & Alto &
  Anonimato estricto, transparencia de criterios y uso exclusivamente
  formativo de los resultados. \\
R-05 & Retraso en la integración con el sistema central. &
  Media & Medio &
  Capa de abstracción (\ac{API} \ac{REST}) y datos simulados para pruebas
  tempranas. \\
R-06 & Bajo desempeño de los modelos analíticos. &
  Media & Medio &
  Validación cruzada, comparación contra línea base y ajuste iterativo
  de variables. \\
R-07 & Incumplimiento de plazos por carga operativa del equipo. &
  Alta & Medio &
  Priorización por valor, entregas incrementales y revisión semanal del
  cronograma. \\
R-08 & Incidente de seguridad o fuga de información. &
  Baja & Muy alto &
  Cifrado, \ac{RBAC}, auditoría de accesos, respaldo y anonimización de
  \ac{PII}. \\
\bottomrule
\end{tabularx}
\end{table}

\section{Recursos Requeridos}

\begin{table}[H]
\centering
\caption{Recursos humanos, técnicos y de infraestructura}
\label{tab:recursos}
\small
\begin{tabularx}{\textwidth}{@{}l X@{}}
\toprule
\rowcolor{verdeClaro}
\textbf{Categoría} & \textbf{Recursos} \\
\midrule
Humanos & Investigador principal, analista de datos, desarrollador \ac{API},
  diseñador de interfaz, asesor metodológico, contraparte de secretaría. \\
Software & Entorno de desarrollo, gestor de base de datos relacional,
  herramienta de analítica, editor de documentación, control de versiones. \\
Infraestructura & Servidor de aplicaciones, servidor de base de datos,
  repositorio analítico, dominio institucional con certificado TLS. \\
Datos & Registros históricos de cupos, materias, faltas, evaluaciones y
  quejas del \ac{DIC}, debidamente anonimizados. \\
\bottomrule
\end{tabularx}
\end{table}

\section{Presupuesto Estimado}

\begin{table}[H]
\centering
\caption{Presupuesto estimado del proyecto}
\label{tab:presupuesto}
\small
\begin{tabularx}{\textwidth}{@{}l X r r@{}}
\toprule
\rowcolor{azulClaro}
\textbf{Rubro} & \textbf{Descripción} & \textbf{Unidades} & \textbf{\%} \\
\midrule
Talento humano & Investigación, análisis de datos y desarrollo & 60 & 60\% \\
Infraestructura & Servidores y almacenamiento (periodo del proyecto) & 15 & 15\% \\
Licencias & Herramientas de analítica y notificaciones & 10 & 10\% \\
Capacitación & Formación de docentes y secretaría & 8 & 8\% \\
Imprevistos & Contingencia del proyecto & 7 & 7\% \\
\midrule
\rowcolor{grisClaro}
\textbf{Total} & & \textbf{100} & \textbf{100\%} \\
\bottomrule
\end{tabularx}
\end{table}

\section{Criterios de Aceptación}

\begin{table}[H]
\centering
\caption{Criterios de aceptación del proyecto}
\label{tab:aceptacion}
\small
\begin{tabularx}{\textwidth}{@{}p{1.4cm} X@{}}
\toprule
\rowcolor{moradoClaro}
\textbf{ID} & \textbf{Criterio verificable} \\
\midrule
CA-01 & Los 56 requerimientos funcionales del Capítulo~\ref{ch:requerimientos}
  están trazados a objetivos estratégicos y casos de uso. \\
CA-02 & Los diagramas \ac{UML} compilan en \LaTeX{} y son consistentes con la
  especificación de requerimientos. \\
CA-03 & El tiempo medio de respuesta a solicitudes de cupo es inferior a 24 horas
  durante la prueba piloto. \\
CA-04 & El 100\% de las evaluaciones docentes del periodo piloto se completan
  dentro de la ventana configurada. \\
CA-05 & El 95\% de las quejas registradas cuentan con responsable, plazo y estado
  de trazabilidad. \\
CA-06 & El modelo de preferencia horaria alcanza una aceptación promedio superior
  a 4.0 sobre 5 en el grupo piloto. \\
CA-07 & El motor analítico genera al menos 10 recomendaciones priorizadas y el
  70\% de ellas son aplicadas por la secretaría o el departamento. \\
CA-08 & El tablero de verificación está disponible para secretaría y
  departamento con todos los paneles definidos. \\
\bottomrule
\end{tabularx}
\end{table}

\section{Conclusiones}

\begin{cajaObjetivo}[Conclusiones del documento]
\begin{enumerate}[leftmargin=1.2cm, itemsep=4pt]
  \item El problema del \ac{DIC} \textbf{no es de capacidad tecnológica sino de
        integración}: los procesos existen, pero operan de manera aislada, sin
        trazabilidad ni analítica, lo que provoca demoras, evaluaciones
        incompletas y decisiones sin evidencia.

  \item Se evidenció que \textbf{once procesos críticos} concentran la mayor
        parte de las dificultades: cupos, materias, comunicación bidireccional,
        faltas, evaluación docente, quejas, efectividad de la secretaría,
        verificación de resultados, asignación de horarios y ausencia de
        analítica.

  \item La metodología \textbf{mixta con diseño de investigación--acción} resulta
        adecuada porque permite comprender las causas humanas del problema y, al
        mismo tiempo, medir objetivamente variables como tiempos de respuesta,
        tasas de participación y frecuencia de quejas.

  \item El uso de \ac{CRISP-DM} ofrece una ruta estructurada y auditable para el
        componente analítico, garantizando que las recomendaciones se evalúen
        contra una línea base antes de desplegarse.

  \item El \textbf{modelado \ac{UML} completo} (casos de uso, clases, secuencia,
        actividades, estados y componentes) constituye una especificación
        suficiente para iniciar la construcción del sistema sin ambigüedades
        funcionales.

  \item El \textbf{modelo de preferencias por probabilidad} convierte los
        cuestionarios de horarios en una herramienta de decisión objetiva,
        reemplazando la programación académica basada en intuición por una
        basada en evidencia.

  \item El \textbf{algoritmo de soluciones rápidas} con priorización por
        impacto, urgencia y esfuerzo permite a la secretaría y al departamento
        actuar sobre lo verdaderamente relevante, y no solo sobre lo más
        visible.

  \item La \textbf{verificación de resultados mediante tableros} cierra el ciclo
        de mejora continua: cada acción queda registrada, medida y disponible
        para la siguiente iteración de decisión.
\end{enumerate}
\end{cajaObjetivo}

\section{Recomendaciones}

\begin{enumerate}[leftmargin=1.4cm, itemsep=4pt]
  \item \textbf{Implementar por incrementos:} iniciar con los módulos de cupos,
        comunicación, faltas y evaluación docente, que concentran el mayor
        impacto sobre la oportunidad.

  \item \textbf{Institucionalizar el porcentaje mínimo de participación:} la
        evaluación docente debe ser requisito para la matrícula del siguiente
        periodo, con anonimato garantizado.

  \item \textbf{Capacitar a la secretaría en el uso analítico:} no basta con
        entregar tableros; se requiere formación en interpretación de
        indicadores y priorización.

  \item \textbf{Publicar los criterios de priorización de cupos:} la
        transparencia es requisito para la legitimidad del sistema ante los
        estudiantes.

  \item \textbf{Revisar periódicamente las reglas del algoritmo:} los umbrales
        (saturación, inasistencia, plazos) deben ajustarse con datos reales de
        cada periodo académico.

  \item \textbf{Diseñar el plan de protección de datos:} clasificar la
        información, definir tiempos de retención y establecer controles de
        acceso antes de la puesta en producción.

  \item \textbf{Extender progresivamente a otras unidades académicas:} la
        arquitectura modular permite replicar el modelo, adaptando reglas de
        negocio y catálogos.

  \item \textbf{Medir el impacto social:} además de los \ac{KPI} operativos, se
        recomienda medir la percepción de equidad, la reducción de conflictos y
        el clima académico percibido.
\end{enumerate}

\section{Trabajo Futuro}

\begin{table}[H]
\centering
\caption{Líneas de trabajo futuro}
\label{tab:trabajo-futuro}
\small
\begin{tabularx}{\textwidth}{@{}p{3.9cm} X c@{}}
\toprule
\rowcolor{azulClaro}
\textbf{Línea} & \textbf{Descripción} & \textbf{Horizonte} \\
\midrule
Asistente conversacional &
  Interfaz de lenguaje natural para consultar cupos, horarios y estados de casos. &
  Corto \\
Recomendación personalizada &
  Sugerencia de materias y horarios según historial y perfil del estudiante. &
  Mediano \\
Detección temprana de deserción &
  Modelos predictivos sobre inasistencia, desempeño y quejas. &
  Mediano \\
Integración total con el sistema central &
  Sincronización bidireccional de matrícula, notas y titulación. &
  Mediano \\
Analítica de aprendizaje &
  Correlación entre desempeño docente y resultados académicos de los estudiantes. &
  Largo \\
Aplicación móvil nativa &
  Notificaciones push y operaciones sin conexión. &
  Largo \\
\bottomrule
\end{tabularx}
\end{table}

% ==================== REFERENCIAS ====================================
\begin{thebibliography}{99}
\addcontentsline{toc}{chapter}{Referencias}

\bibitem{omg-uml}
Object Management Group. \textit{OMG Unified Modeling Language (OMG UML), Version 2.5.1}. OMG, 2017.

\bibitem{sommerville}
Sommerville, I. \textit{Software Engineering}. 10th ed. Pearson, 2016.

\bibitem{pressman}
Pressman, R. S. y Maxim, B. R. \textit{Software Engineering: A Practitioner's Approach}. 9th ed. McGraw-Hill, 2020.

\bibitem{pmbok}
Project Management Institute. \textit{A Guide to the Project Management Body of Knowledge (PMBOK Guide)}. 7th ed. PMI, 2021.

\bibitem{crispdm}
Chapman, P., Clinton, J., Kerber, R., Khabaza, T., Reinartz, T., Shearer, C. y Wirth, R. \textit{CRISP-DM 1.0: Step-by-step Data Mining Guide}. SPSS Inc., 2000.

\bibitem{iso25010}
International Organization for Standardization. \textit{ISO/IEC 25010: Systems and software engineering --- Systems and software Quality Requirements and Evaluation (SQuaRE)}. ISO, 2011.

\bibitem{iso27001}
International Organization for Standardization. \textit{ISO/IEC 27001: Information security, cybersecurity and privacy protection --- Information security management systems}. ISO, 2022.

\end{thebibliography}

\end{document}
